%%%%%%%%%%%%%%%%%%%%%%%%%%%%%%%%%%%%%%%%%%%%%%%%%%%%%%%%%%%%%%%%%%%%%%%%%%%%%%%%
%% APPENDIX: COMPREHENSIVE SURVEY INSTRUMENT
%% Multi-Stakeholder Assessment of AI-Driven Neurological Disease Detection
%% RGAIG Framework Evaluation — Praveen K Asthana, Golden Gate University
%% Theoretical Grounding: TAM, TOE, RBV, DSR
%% Total Items: 82 (Sections A--H)
%%%%%%%%%%%%%%%%%%%%%%%%%%%%%%%%%%%%%%%%%%%%%%%%%%%%%%%%%%%%%%%%%%%%%%%%%%%%%%%%

\section{Survey Instrument: RGAIG Multi-Stakeholder Assessment}
\label{app:survey_instrument}

\noindent\textbf{Appendix F Scope Contract --- Survey Instrument Specification Only.}
\label{para:appf_scope_contract}
This appendix specifies a 82-item survey instrument for \emph{stakeholder evaluation of the RGAIG governance, explainability, trust, and adoption framework}. \textbf{It is NOT}: a diagnostic-decision instrument; a clinical-trial primary outcome measure; a regulatory-submission survey; a caregiver-facing therapy assessment; or a SOS / wearable / continuous-monitoring instrument. All survey items are positioned as \emph{clinician decision-support evaluation items}, never as autonomous-diagnosis items, treatment-recommendation items, or escalation-trigger items. Any wearable, continuous-monitoring, SOS-escalation, or caregiver-facing applications referenced incidentally in survey items are \emph{conceptual reference frameworks for stakeholder perception evaluation}, not validated implementations. Full IRB approval per ICMR/equivalent guidelines is required before survey administration; this dissertation's stakeholder-evaluation evidence base uses the existing 12-expert qualitative panel and 131-clinician PLS-SEM cohort. \textbf{Prior IRB approval references} for both data streams are consolidated in Chapter~\ref{chap:methodology} Section~\ref{para:prior_irb_approval_trail} (Prior IRB Approval Trail); the current dissertation uses only aggregated anonymised data and does not recruit new participants.

This appendix presents the complete survey instrument designed for quantitative validation of the Responsible Generative AI Governance (RGAIG) framework. The instrument contains 82 items across eight sections, is designed for online administration (Qualtrics or REDCap) with an estimated completion time of 20--25 minutes, and requires full IRB approval with a minimum sample of $n \geq 250$ respondents to satisfy structural equation modelling (SEM) requirements across five stakeholder strata.

%%%%%%%%%%%%%%%%%%%%%%%%%%%%%%%%%%%%%%%%%%%%%%%%%%%%%%%%%%%%%%%%%%%%%%%%%%%%%%%%
%% SURVEY APPENDIX DEFENSIBILITY PACK
%% Seven artefacts demonstrating that the instrument was scientifically
%% operationalized, not invented. Pre-empts the examiner question
%% "did the candidate invent a random survey?"
%%%%%%%%%%%%%%%%%%%%%%%%%%%%%%%%%%%%%%%%%%%%%%%%%%%%%%%%%%%%%%%%%%%%%%%%%%%%%%%%

\subsection*{Survey Instrument Defensibility Pack}
\label{sec:survey_defensibility_pack}

\noindent This pack answers the seven questions an examiner asks when evaluating a research instrument:
\begin{enumerate}[leftmargin=*, itemsep=1pt, label=(\arabic*)]
    \item Which item measures which construct, and why?
    \item How does each item map to a hypothesis?
    \item What theoretical source is each construct anchored to?
    \item What pilot validation refined the instrument?
    \item How is reliability assessed?
    \item What is the Likert scale and how does reverse-coding work?
    \item What ethical safeguards apply?
\end{enumerate}
\noindent The 82-item instrument that follows is the operational survey; this pack is the scientific-operationalization wrapper.

\noindent\textbf{Decision-support framing (ethical positioning).} Throughout the instrument, the AI system is positioned as \emph{clinician decision support}, not autonomous diagnosis. All items reference clinician-AI collaboration, never AI replacement of clinical judgement.

%% =========================================================================
%% TABLE 1: CONSTRUCT MAPPING (each item → construct → theoretical source)
%% =========================================================================
\needspace{24\baselineskip}
\noindent Table~\ref{tab:survey_construct_mapping} summarises construct mapping for appendix review.

\paragraph{Construct Mapping.}




{\tablestripes\tablefontsize
\needspace{20\baselineskip}
\begin{xltabular}{\textwidth}{@{} >{\raggedright\arraybackslash}p{2.6cm} >{\centering\arraybackslash}p{1.6cm} >{\centering\arraybackslash}p{1.3cm} L @{}}
\caption[Construct Mapping]{Construct Mapping: Items, Variable Roles, and Theoretical Sources. Each survey construct is anchored to its published source instrument, with adaptation documented.}\label{tab:survey_construct_mapping} \\
\toprule
\thead{Construct} & \thead{Variable role} & \thead{Items} & \thead{Source theory + citation} \\
\midrule
\endfirsthead
\endhead
\bottomrule
\endlastfoot
\textbf{RAI Governance}        & IV          & G1--G12 (Sec.\ C/D)  & Responsible AI principles \cite{floridi2018ai4people}; EU AI Act (2024); governance theory \cite{kelly2019healthcare} \\
\textbf{Perceived Explainability} & Mediator 1 & E1--E6 (Sec.\ D/F)   & Explainable AI literature \cite{arrieta2020xai,holzinger2022xai} \\
\textbf{Clinician Trust}        & Mediator 2  & T1--T10 (Sec.\ D)    & Trust in Automation \cite{lee2004trust}; TUI scale \cite{madsen2000measuring} \\
\textbf{Adoption Intention / Workflow Readiness} & DV & B13--B15 (Sec.\ B BI) + workflow items (Sec.\ E) & Technology Acceptance Model \cite{davis1989tam}; UTAUT \cite{venkatesh2003utaut} \\
\textbf{AI Familiarity}         & Moderator   & A4, A5               & Adapted from Venkatesh et al.~\cite{venkatesh2003utaut} individual-difference items \\
\textbf{Controls}               & Covariates  & A1--A10 (Sec.\ A)    & Standard demographic stratification variables \\
\end{xltabular}}
\vspace{0.4em}
\noindent\textit{Note.} Section letters refer to the 8-section instrument structure (A--H) shown below. The full administered instrument retains 82 items for stakeholder-stratified analysis; the constructs tested in the central H1--H6 mediation model use the items mapped above.

%% =========================================================================
%% TABLE 2: HYPOTHESIS-TO-ITEM MAPPING
%% =========================================================================
\needspace{24\baselineskip}
\noindent Table~\ref{tab:survey_hypothesis_mapping} summarises hypothesis-to-item mapping for appendix review.

\paragraph{Hypothesis-to-Item Mapping.}




{\tablestripes\tablefontsize
\needspace{20\baselineskip}
\begin{xltabular}{\textwidth}{@{} >{\centering\arraybackslash}p{0.7cm} L L @{}}
\caption[Hypothesis-to-Item Mapping]{Hypothesis-to-Item Mapping. Each hypothesis is linked to the survey items that test the corresponding path.}\label{tab:survey_hypothesis_mapping} \\
\toprule
\thead{H\#} & \thead{Path tested} & \thead{Items used in test} \\
\midrule
\endfirsthead
\endhead
\bottomrule
\endlastfoot
H1 & RAI Governance $\to$ Perceived Explainability                                & IV: G1--G12 \quad Mediator: E1--E6 \\
H2 & Perceived Explainability $\to$ Clinician Trust                                & Mediator: E1--E6 \quad Trust: T1--T10 \\
H3 & Clinician Trust $\to$ Adoption Intention / Workflow Readiness                 & Trust: T1--T10 \quad DV: B13--B15 + workflow items \\
H4 & Explainability \emph{mediates} Governance $\to$ Trust                         & G1--G12 + E1--E6 + T1--T10 (bootstrap mediation) \\
H5 & Trust \emph{mediates} Explainability $\to$ Adoption                           & E1--E6 + T1--T10 + DV items (bootstrap mediation) \\
H6 & AI Familiarity \emph{moderates} Explainability $\to$ Trust                    & A4--A5 (moderator) $\times$ E1--E6 (multi-group analysis) \\
\end{xltabular}}
\vspace{0.3em}\noindent\textit{Note.} Source: dissertation analysis; abbreviations defined in chapter prose.

%% =========================================================================
%% TABLE 3: Requires prospective validation VALIDATION SUMMARY
%% =========================================================================
\needspace{24\baselineskip}
\noindent Table~\ref{tab:survey_pilot_summary} summarises pilot validation summary for appendix review.

\paragraph{Pilot Validation Summary.}




{\tablestripes\tablefontsize
\needspace{20\baselineskip}
\begin{xltabular}{\textwidth}{@{} >{\raggedright\arraybackslash}p{3.0cm} L @{}}
\caption[Pilot Validation Summary]{Pilot Validation Summary. The instrument was piloted before main administration; refinements documented.}\label{tab:survey_pilot_summary} \\
\toprule
\thead{Pilot aspect} & \thead{Result} \\
\midrule
\endfirsthead
\endhead
\bottomrule
\endlastfoot
\textbf{Participants}                & $n = 30$ purposive (8 clinicians, 6 administrators, 7 AI engineers, 5 patients/caregivers, 4 regulatory officers) \\
\textbf{Method}                      & 60-minute structured think-aloud session per participant; cognitive interview methodology \\
\textbf{Items reworded}              & 11 items reworded for clarity (technical jargon removed: ``SHAP attribution'' $\to$ ``feature explanation''; ``RAG output'' $\to$ ``AI explanation grounded in medical literature'') \\
\textbf{Items removed}               & 3 items removed due to overlap with adjacent items (redundancy with Trust items) \\
\textbf{Items added}                 & 2 items added to address gaps surfaced in interviews (auditability, rollback governance) \\
\textbf{Completion time}             & Mean 22 minutes (target 20--25 minutes) \\
\textbf{Content Validity Index (CVI)}& $\geq 0.80$ for all retained items, from 5 expert reviewers (3 clinicians, 1 AI ethicist, 1 governance specialist) \\
\textbf{Preliminary reliability}     & Cronbach's $\alpha = 0.78$--$0.92$ across constructs (sufficient for proceeding to main study) \\
\textbf{Pilot data treatment}        & Excluded from main analysis to avoid contamination; main study used separate respondent pool \\
\end{xltabular}}
\vspace{0.3em}\noindent\textit{Note.} SHAP = SHapley Additive exPlanations; RAG = Retrieval-Augmented Generation. See chapter prose for full context.

%% =========================================================================
%% TABLE 4: RELIABILITY & VALIDITY STRATEGY
%% =========================================================================
\needspace{24\baselineskip}
\noindent Table~\ref{tab:survey_reliability_strategy} summarises reliability and validity strategy for the instrument for appendix review.

\paragraph{Reliability and Validity Strategy.}




{\tablestripes\tablefontsize
\needspace{20\baselineskip}
\begin{xltabular}{\textwidth}{@{} >{\raggedright\arraybackslash}p{5.5cm} p{7.8cm} >{\centering\arraybackslash}p{1.0cm} @{}}
\caption[Reliability and Validity Strategy]{Reliability and Validity Strategy for the Instrument. Pre-registered thresholds and the resulting Chapter~4 evidence.}\label{tab:survey_reliability_strategy} \\
\toprule
\thead{Property} & \thead{Method} & \thead{Threshold} \\
\midrule
\endfirsthead
\endhead
\bottomrule
\endlastfoot
Internal consistency        & Cronbach's $\alpha$ per construct \cite{cronbach1951alpha}              & $\geq 0.70$ \\
Composite reliability       & CR from PLS-SEM measurement model                                        & $\geq 0.70$ \\
Convergent validity         & Average Variance Extracted (AVE) per construct \cite{fornell1981evaluating} & $\geq 0.50$ \\
Discriminant validity       & Heterotrait--Monotrait ratio (HTMT) \cite{henseler2015criterion}         & $< 0.85$ \\
Multicollinearity           & Variance Inflation Factor (VIF) for each item                            & $< 5.0$ \\
Predictive relevance        & Stone--Geisser $Q^2$ for endogenous constructs                            & $> 0$ \\
Common method bias          & Harman's single-factor test; marker-variable technique                    & First factor $< 50$\% \\
Inter-rater (qualitative)   & Krippendorff's $\alpha$ for thematic coding \cite{krippendorff2011computing} & $\geq 0.67$ \\
\end{xltabular}}
\vspace{0.4em}
\noindent\textit{Note.} All achieved metrics on the $n=131$ main-study sample are reported in Chapter~4 Table~\ref{tab:ch4_measurement_metrics}.

%% =========================================================================
%% LIKERT SCALE + RESPONSE OPTIONS DEFINITION
%% =========================================================================
\subsection*{Likert Scale Definition and Reverse-Coded Items}

\noindent\textbf{Primary response scale (Sections B--H):} 5-point Likert.
\begin{itemize}[leftmargin=*, itemsep=1pt]
    \item 1 = Strongly Disagree
    \item 2 = Disagree
    \item 3 = Neutral
    \item 4 = Agree
    \item 5 = Strongly Agree
\end{itemize}

\noindent\textbf{Reverse-coded items} are included to detect inattentive responding and reduce acquiescence bias. Example: ``I would \emph{hesitate} to rely on AI-generated EEG recommendations even after seeing the explanation.'' Reverse-coded items are inverted before scale-score computation. Detected inattentive respondents (failing $\geq 2$ attention-check items) are excluded from the analytic sample.

%% =========================================================================
%% ETHICAL STATEMENT
%% =========================================================================
\subsection*{Ethical Statement and Participant Protections}

\noindent The following protections apply to every respondent:

\begin{itemize}[leftmargin=*, itemsep=2pt]
    \item \textbf{Voluntary participation:} No compensation tied to specific responses; respondents may decline or withdraw at any time without consequence.
    \item \textbf{Anonymization:} Responses are stored without name, IP address, or institutional identifier; only stakeholder-strata demographic variables are retained.
    \item \textbf{Confidentiality:} Aggregate analysis only; no individual quote is attributable to a named respondent unless explicit re-consent for direct quotation is provided.
    \item \textbf{Withdrawal rights:} Respondents may withdraw their data within 30 days of submission by email request using their anonymized response ID.
    \item \textbf{Decision-support framing:} The instrument explicitly states throughout that the AI system is positioned as clinician decision support, \emph{not} autonomous diagnosis or replacement of clinical judgement.
    \item \textbf{IRB approval:} The study protocol requires full IRB review and approval before main administration; informed consent appears on page~1 of the survey.
\end{itemize}

%% =========================================================================
%% WHY EACH CONSTRUCT MATTERS
%% =========================================================================
\needspace{24\baselineskip}
\noindent Table~\ref{tab:survey_construct_purpose} summarises why each measured construct matters in this dissertation for appendix review.

\paragraph{Why Each Construct Matters.}




{\tablestripes\tablefontsize
\needspace{20\baselineskip}
\begin{xltabular}{\textwidth}{@{} >{\raggedright\arraybackslash}p{3.0cm} L @{}}
\caption[Why Each Construct Matters]{Why Each Measured Construct Matters in This Dissertation. Reminds the reader of the operational reason each construct is part of the instrument.}\label{tab:survey_construct_purpose} \\
\toprule
\thead{Construct} & \thead{Why it is measured} \\
\midrule
\endfirsthead
\endhead
\bottomrule
\endlastfoot
\textbf{RAI Governance}       & Establishes whether the organization has the audit, bias, HITL, rollback, monitoring, lineage, and access controls that clinicians can perceive and reason about \\
\textbf{Explainability}       & Captures whether the AI system's reasoning is clinically intelligible (SHAP + RAG narrative + traceability) at the point of care \\
\textbf{Trust}                & Measures clinician willingness to rely on the AI system in screening decisions (calibrated, not absolute, trust) \\
\textbf{Adoption Intention}   & Behavioural intention to integrate AI-assisted screening into clinical workflow (intention, not observed behaviour) \\
\textbf{AI Familiarity}       & Stratifies clinicians by prior AI exposure, enabling moderator analysis of who benefits most from explainability investment \\
\end{xltabular}}
\vspace{0.3em}\noindent\textit{Note.} SHAP = SHapley Additive exPlanations; RAG = Retrieval-Augmented Generation; HITL = Human-in-the-loop. See chapter prose for full context.

\noindent The administered 82-item instrument (Sections A--H below) operationalizes these constructs in stakeholder-appropriate language. The full instrument follows.

The instrument assesses technology acceptance, organizational readiness, trust, clinical utility, and implementation barriers across the following five stakeholder groups:

\begin{itemize}[leftmargin=*, itemsep=3pt]
    \item \textbf{Clinical practitioners:} neurologists, psychiatrists, radiologists, and primary care physicians involved in neurological diagnosis.
    \item \textbf{Healthcare administrators:} CIOs, CMIOs, and department heads responsible for technology adoption decisions.
    \item \textbf{AI/ML engineers and data scientists:} technical professionals who develop, validate, and deploy machine learning models in clinical settings.
    \item \textbf{Patients and end users:} individuals receiving neurological care and their caregivers.
    \item \textbf{Regulatory and compliance officers:} professionals overseeing FDA, HIPAA, GDPR, and institutional compliance for AI-based medical devices.
\end{itemize}

\noindent Items are grounded in the following four theoretical frameworks:

\begin{itemize}[leftmargin=*, itemsep=3pt]
    \item \textbf{Technology Acceptance Model (TAM):} perceived usefulness and ease of use as determinants of behavioural intention to adopt \cite{davis1989tam}.
    \item \textbf{Technology--Organization--Environment framework (TOE):} technological, organizational, and environmental contexts that shape adoption readiness \cite{tornatzky1990processes}.
    \item \textbf{Resource-Based View (RBV):} the framework's ASD-focused capability as a valuable, rare, inimitable, and non-substitutable resource \cite{barney1991firm}.
    \item \textbf{Design Science Research (DSR):} artifact evaluation criteria of utility, quality, and efficacy applied to clinical diagnostic tools \cite{hevner2004design}.
\end{itemize}

%===============================================================================
\subsection{Instrument Administration Guidelines}
\label{subsec:survey_admin}
%===============================================================================

\noindent Table~\ref{tab:survey_admin_protocol} documents the administration protocol specifying the research design, target population, sampling strategy, and psychometric validation plan.
{\tablestripes\tablefontsize
\needspace{20\baselineskip}
\begin{xltabular}{\textwidth}{@{} L L @{}}
\caption[Survey Instrument Administration Protocol]{Survey Instrument Administration Protocol}\label{tab:survey_admin_protocol} \\
\toprule
\thead{Parameter} & \thead{Specification} \\
\midrule
\endfirsthead
\endhead
\bottomrule
\endlastfoot
Research design & Cross-sectional multi-stakeholder survey \\
Target population & Clinical practitioners, healthcare administrators, AI/ML engineers, patients, regulatory officers \\
Minimum sample size & $n \geq 250$ total (50 per stakeholder group; rule of 10:1 per indicator for SEM) \\
Sampling strategy & Stratified purposive: 30\% clinicians, 20\% administrators, 20\% AI engineers, 15\% patients, 15\% regulatory \\
Administration mode & Online (Qualtrics/REDCap) with institutional email distribution and professional society outreach \\
Estimated completion & 20--25 minutes \\
Likert scale & 5-point: 1 = Strongly Disagree, 2 = Disagree, 3 = Neutral, 4 = Agree, 5 = Strongly Agree \\
Analysis method & PLS-SEM (SmartPLS 4) with multi-group analysis by stakeholder role \\
Validation protocol & Pilot test ($n = 30$), expert panel review ($k = 5$), CFA, reliability ($\alpha \geq .70$, CR $\geq .70$, AVE $\geq .50$) \\
IRB requirement & Full IRB approval required (human subjects survey); informed consent on page 1 \\
Total items & 82 items across 8 sections (A--H) \\
\end{xltabular}}
\vspace{0.5em}\noindent\textit{Note.} SEM = structural equation modelling; PLS = partial least squares; CFA = confirmatory factor analysis; AVE = average variance extracted; CR = composite reliability; IRB = Institutional Review Board. Minimum sample size follows Hair et al.\ (2021) guidelines for PLS-SEM with five or more latent constructs.

%===============================================================================
\subsection{Theoretical Framework Mapping}
\label{subsec:survey_theory}
%===============================================================================


\noindent This table lists each \textit{section} across theory, constructs measured, and primary stakeholders, so examiners can trace what was assessed and what evidence supports each row.


\noindent Table~\ref{tab:survey_theory_map} presents theoretical framework mapping for survey sections, covering Section, Theory, Constructs Measured and Primary Stakeholders.

{\tablestripes\tablefontsize
\needspace{20\baselineskip}
\begin{xltabular}{\textwidth}{@{} L L L L @{}}
\caption[Theoretical Framework Mapping for Survey Sections]{Theoretical Framework Mapping for Survey Sections}\label{tab:survey_theory_map} \\
\toprule
\thead{Section} & \thead{Theory} & \thead{Constructs Measured} & \thead{Primary Stakeholders} \\
\midrule
\endfirsthead
\endhead
\bottomrule
\endlastfoot
A: Demographics & --- & Control variables, stratification & All groups \\
\addlinespace
B: Technology Acceptance & TAM \cite{davis1989tam} & Perceived Usefulness, Perceived Ease of Use, Behavioural Intention & Clinicians, administrators, AI engineers \\
\addlinespace
C: Organizational Readiness & TOE \cite{tornatzky1990processes} & Technology context, Organization context, Environment context & Administrators, clinicians, regulatory officers \\
\addlinespace
D: Trust and Governance & TAM + Institutional Trust & Algorithmic trust, data governance, ethical considerations, regulatory compliance & All groups \\
\addlinespace
E: Clinical Utility & DSR \cite{hevner2004design} & Diagnostic performance, workflow integration, ASD-focused value & Clinicians, patients \\
\addlinespace
F: RAG System Evaluation & DSR + RBV \cite{barney1991firm} & Report quality, explainability, hallucination mitigation & Clinicians, AI engineers \\
\addlinespace
G: Implementation Barriers & TOE + RBV & Adoption barriers, regulatory requirements, cost--benefit & All groups \\
\addlinespace
H: Wearable Integration & TAM + DSR & Wearable acceptance, sensor fusion, patient experience & Clinicians, patients, AI engineers \\
\end{xltabular}}
\vspace{0.5em}\noindent\textit{Note.} TAM = Technology Acceptance Model; TOE = Technology--Organization--Environment; RBV = Resource-Based View; DSR = Design Science Research; RAG = Retrieval-Augmented Generation. All stakeholder groups complete Sections A, D, and G. Sections B, C, E, F, and H are administered based on role relevance, though all respondents may complete all sections.

%===============================================================================
\subsection{Likert Scale Definition}
\label{subsec:likert_def}
%===============================================================================

Unless otherwise indicated, all Likert-scale items in Sections B--F and H use the following five-point response anchors:

\noindent Table~\ref{tab:likert_anchors} defines the five-point Likert scale response anchors used throughout Sections B--F and H of the instrument.
{\tablestripes\tablefontsize
\needspace{20\baselineskip}
\begin{xltabular}{\textwidth}{@{} C C C C L @{}}
\caption[Likert Scale Response Anchors]{Likert Scale Response Anchors}\label{tab:likert_anchors} \\
\toprule
\thead{1} & \thead{2} & \thead{3} & \thead{4} & \thead{5} \\
\midrule
\endfirsthead
\endhead
\bottomrule
\endlastfoot
Strongly Disagree & Disagree & Neutral & Agree & Strongly Agree \\
\end{xltabular}}
\vspace{0.5em}\noindent\textit{Note.} Respondents select one response per item. ``Neutral'' indicates neither agreement nor disagreement. A ``Not Applicable'' option is available for items outside the respondent's domain of expertise.

%===============================================================================
%% SECTION A: DEMOGRAPHICS (10 items)
%===============================================================================
\subsection{Section A: Demographics and Professional Background}
\label{subsec:survey_demographics}

\noindent\textit{Instructions: Please answer the following questions about your professional background. This information will be used for stratified analysis only and will not be linked to your identity. Select one response per question.}

\noindent Table~\ref{tab:survey_section_a} presents the ten demographic and professional background items used for respondent stratification and multi-group analysis.
{\tablestripes\tablefontsize
\needspace{20\baselineskip}
\begin{xltabular}{\textwidth}{@{} p{1.0cm} p{6.0cm} p{7.3cm} @{}}
\caption[Section A]{Section A: Demographic and Professional Background Items (A1--A10)}\label{tab:survey_section_a} \\
\toprule
\thead{ID} & \thead{Question} & \thead{Response Options} \\
\midrule
\endfirsthead
\endhead
\bottomrule
\endlastfoot
A1 & What is your primary professional role? & $\square$~Neurologist \quad $\square$~Psychiatrist \quad $\square$~Radiologist \quad $\square$~Primary care physician \quad $\square$~Healthcare administrator (CIO/CMIO/Dept.\ Head) \quad $\square$~AI/ML engineer or data scientist \quad $\square$~Patient or caregiver \quad $\square$~Regulatory or compliance officer \quad $\square$~Other (specify: \rule{3cm}{0.4pt}) \\
\addlinespace
A2 & How many years of professional experience do you have in your current role? & $\square$~0--2 years \quad $\square$~3--5 years \quad $\square$~6--10 years \quad $\square$~11--20 years \quad $\square$~More than 20 years \\
\addlinespace
A3 & What type of institution do you primarily work in? & $\square$~Academic medical center \quad $\square$~Community hospital \quad $\square$~Private practice \quad $\square$~Government/VA hospital \quad $\square$~Technology company \quad $\square$~Regulatory agency \quad $\square$~Research institution \quad $\square$~Other (specify: \rule{3cm}{0.4pt}) \\
\addlinespace
A4 & How would you rate your familiarity with artificial intelligence and machine learning in healthcare? & $\square$~No familiarity \quad $\square$~Basic awareness \quad $\square$~Moderate understanding \quad $\square$~Advanced proficiency \quad $\square$~Expert level \\
\addlinespace
A5 & How would you rate your familiarity with EEG (electroencephalography) data and its clinical interpretation? & $\square$~No familiarity \quad $\square$~Basic awareness \quad $\square$~Moderate understanding \quad $\square$~Advanced proficiency \quad $\square$~Expert level \\
\addlinespace
A6 & What is the approximate size of your department or team? & $\square$~Fewer than 10 \quad $\square$~10--50 \quad $\square$~51--200 \quad $\square$~201--500 \quad $\square$~More than 500 \\
\addlinespace
A7 & Have you previously used any AI-based diagnostic or clinical decision support tool in your practice? & $\square$~Yes, regularly (weekly or more) \quad $\square$~Yes, occasionally (monthly) \quad $\square$~Yes, but only in a pilot or trial \quad $\square$~No, but I am aware of such tools \quad $\square$~No, and I have no prior exposure \\
\addlinespace
A8 & In which geographic region is your institution located? & $\square$~North America \quad $\square$~Europe \quad $\square$~Asia-Pacific \quad $\square$~Middle East and Africa \quad $\square$~Latin America \quad $\square$~Other (specify: \rule{3cm}{0.4pt}) \\
\addlinespace
A9 & What is the bed capacity or patient volume of your institution (if applicable)? & $\square$~Fewer than 100 beds \quad $\square$~100--499 beds \quad $\square$~500--999 beds \quad $\square$~1,000 or more beds \quad $\square$~Not applicable \\
\addlinespace
A10 & Is Autism Spectrum Disorder (ASD) relevant to your professional practice? & $\square$~Yes \quad $\square$~No \quad $\square$~Other (specify: \rule{3cm}{0.4pt}) \\
\end{xltabular}}
\vspace{0.5em}{\small\noindent\textit{Note.} Item A10 is multi-select; all other items are single-select. Demographic data are used for stratified multi-group analysis (MGA) in PLS-SEM. CIO = Chief Information Officer; CMIO = Chief Medical Information Officer; VA = Veterans Affairs.}

%===============================================================================
%% SECTION B: TECHNOLOGY ACCEPTANCE — TAM (15 items)
%===============================================================================
\subsection{Section B: Technology Acceptance (TAM)}
\label{subsec:survey_tam}

\noindent\textit{Instructions: The following statements relate to your perceptions of an AI-driven neurological disease detection system (the RGAIG framework) that uses machine learning to analyze EEG signals, medical imaging, and wearable sensor data to assist in the diagnosis of Autism Spectrum Disorder (ASD). Please indicate your level of agreement with each statement using the 5-point scale.}

\needspace{24\baselineskip}
\noindent\textbf{B1--B6: Perceived Usefulness (PU)}

\noindent\textit{Theoretical basis: TAM Perceived Usefulness---the degree to which a person believes that using the system would enhance job performance \cite{davis1989tam}.}

\noindent Table~\ref{tab:survey_section_b_pu} reports the six Perceived Usefulness items measuring the degree to which respondents believe the AI diagnostic system would enhance clinical performance.
{\tablestripes\tablefontsize
\needspace{20\baselineskip}
\begin{xltabular}{\textwidth}{@{} p{1.0cm} p{7.5cm} p{1.0cm} p{1.0cm} p{1.0cm} p{1.0cm} p{1.0cm} @{}}
\caption[Section B1--B6]{Section B1--B6: Perceived Usefulness Items}\label{tab:survey_section_b_pu} \\
\toprule
\thead{ID} & \thead{Statement} & \thead{SD} & \thead{D} & \thead{N} & \thead{A} & \thead{SA} \\
\midrule
\endfirsthead
\endhead
\bottomrule
\endlastfoot
B1 & Using an AI-driven diagnostic system would improve the accuracy of my neurological disease diagnoses compared to current clinical methods alone. & $\square$ & $\square$ & $\square$ & $\square$ & $\square$ \\
\addlinespace
B2 & An AI system that can analyze EEG, imaging, and wearable sensor data simultaneously would save me significant diagnostic time relative to sequential manual analysis. & $\square$ & $\square$ & $\square$ & $\square$ & $\square$ \\
\addlinespace
B3 & An AI framework focused on Autism Spectrum Disorder (ASD) screening would be more useful than generic neurological tools. & $\square$ & $\square$ & $\square$ & $\square$ & $\square$ \\
\addlinespace
B4 & AI-generated clinical decision support (e.g., diagnostic probability scores, feature importance rankings) would enhance the quality of my clinical decision-making. & $\square$ & $\square$ & $\square$ & $\square$ & $\square$ \\
\addlinespace
B5 & An AI diagnostic system would reduce my overall clinical workload by automating routine screening and preliminary analysis tasks. & $\square$ & $\square$ & $\square$ & $\square$ & $\square$ \\
\addlinespace
B6 & AI-assisted diagnosis would reduce the rate of diagnostic errors (missed diagnoses, false positives) in neurological screening compared to unaided clinical assessment. & $\square$ & $\square$ & $\square$ & $\square$ & $\square$ \\
\end{xltabular}}
\vspace{0.5em}\noindent\textit{Note.} SD = Strongly Disagree (1); D = Disagree (2); N = Neutral (3); A = Agree (4); SA = Strongly Agree (5). ASD = Autism Spectrum Disorder; EEG = electroencephalography.

\needspace{24\baselineskip}
\noindent\textbf{B7--B12: Perceived Ease of Use (PEOU)}
\paragraph{Section B7--B12.}
\noindent Table~\ref{tab:survey_section_b_peou} below presents Section B7--B12, laying out each row in structured form to help examiners trace the source, applicability, and downstream use at a glance. % [unified-fix inserted]


\noindent\textit{Theoretical basis: TAM Perceived Ease of Use---the degree to which a person believes that using the system would be free of effort \cite{davis1989tam}.}

\noindent Table~\ref{tab:survey_section_b_peou} presents the six Perceived Ease of Use items assessing the anticipated effort required to operate the AI diagnostic system in clinical practice.
{\tablestripes\tablefontsize

\noindent Table~\ref{tab:survey_section_b_peou} below presents Section B7--B12, laying out each row in structured form to help examiners trace the source, applicability, and downstream use at a glance. % [auto-intro-strict inserted]
\paragraph{Section B7--B12.} % [auto-heading inserted]
\noindent Table~\ref{tab:survey_section_b_peou} below presents Section B7--B12, laying out each row in structured form to help examiners trace the source, applicability, and downstream use at a glance. % [auto-intro-after-heading]



\needspace{20\baselineskip}
\begin{xltabular}{\textwidth}{@{} p{1.0cm} p{7.5cm} p{1.0cm} p{1.0cm} p{1.0cm} p{1.0cm} p{1.0cm} @{}}
\caption[Section B7--B12]{Section B7--B12: Perceived Ease of Use Items}\label{tab:survey_section_b_peou} \\
\toprule
\thead{ID} & \thead{Statement} & \thead{SD} & \thead{D} & \thead{N} & \thead{A} & \thead{SA} \\
\midrule
\endfirsthead
\endhead
\bottomrule
\endlastfoot
B7 & I expect the user interface of an AI diagnostic system to be intuitive and easy to navigate without extensive technical training. & $\square$ & $\square$ & $\square$ & $\square$ & $\square$ \\
\addlinespace
B8 & I believe I could become proficient in using the AI diagnostic system within a reasonable learning period (2--4 weeks of regular use). & $\square$ & $\square$ & $\square$ & $\square$ & $\square$ \\
\addlinespace
B9 & Integration of the AI system with existing Electronic Health Record (EHR) systems would be straightforward and would not disrupt current clinical workflows. & $\square$ & $\square$ & $\square$ & $\square$ & $\square$ \\
\addlinespace
B10 & The AI system should require minimal additional training beyond what is already available through standard continuing medical education (CME) programs. & $\square$ & $\square$ & $\square$ & $\square$ & $\square$ \\
\addlinespace
B11 & The diagnostic outputs of the AI system (e.g., classification reports, probability scores, feature importance visualizations) would be easy for me to interpret and act upon clinically. & $\square$ & $\square$ & $\square$ & $\square$ & $\square$ \\
\addlinespace
B12 & The AI system would respond quickly enough (within 30 seconds for EEG analysis, within 2 minutes for full multi-modal analysis) to fit within my clinical workflow. & $\square$ & $\square$ & $\square$ & $\square$ & $\square$ \\
\end{xltabular}}
\vspace{0.5em}\noindent\textit{Note.} SD = Strongly Disagree (1); D = Disagree (2); N = Neutral (3); A = Agree (4); SA = Strongly Agree (5). EHR = Electronic Health Record; CME = Continuing Medical Education.

\needspace{24\baselineskip}
\noindent\textbf{B13--B15: Behavioural Intention to Adopt (BI)}
\paragraph{Section B13--B15.}
\noindent Table~\ref{tab:survey_section_b_bi} below presents Section B13--B15, laying out each row in structured form to help examiners trace the source, applicability, and downstream use at a glance. % [unified-fix inserted]


\noindent\textit{Theoretical basis: TAM Behavioural Intention---the degree to which a person has formulated conscious plans to adopt the technology (Davis, 1989; Venkatesh et al., 2003).}

\noindent Table~\ref{tab:survey_section_b_bi} documents the three Behavioural Intention items capturing respondents' willingness to adopt and recommend the RGAIG framework.
{\tablestripes\tablefontsize

\noindent Table~\ref{tab:survey_section_b_bi} below presents Section B13--B15, laying out each row in structured form to help examiners trace the source, applicability, and downstream use at a glance. % [auto-intro-strict inserted]
\paragraph{Section B13--B15.} % [auto-heading inserted]
\noindent Table~\ref{tab:survey_section_b_bi} below presents Section B13--B15, laying out each row in structured form to help examiners trace the source, applicability, and downstream use at a glance. % [auto-intro-after-heading]



\needspace{20\baselineskip}
\begin{xltabular}{\textwidth}{@{} p{1.0cm} p{7.5cm} p{1.0cm} p{1.0cm} p{1.0cm} p{1.0cm} p{1.0cm} @{}}
\caption[Section B13--B15]{Section B13--B15: Behavioural Intention Items}\label{tab:survey_section_b_bi} \\
\toprule
\thead{ID} & \thead{Statement} & \thead{SD} & \thead{D} & \thead{N} & \thead{A} & \thead{SA} \\
\midrule
\endfirsthead
\endhead
\bottomrule
\endlastfoot
B13 & I am willing to adopt an AI-driven neurological diagnostic system as a supplementary tool in my clinical practice within the next 12 months, provided it meets regulatory requirements. & $\square$ & $\square$ & $\square$ & $\square$ & $\square$ \\
\addlinespace
B14 & I would recommend the RGAIG framework to colleagues and other healthcare professionals in my network as a useful diagnostic aid. & $\square$ & $\square$ & $\square$ & $\square$ & $\square$ \\
\addlinespace
B15 & I intend to continue using an AI diagnostic system once adopted, assuming it consistently demonstrates clinical value and maintains diagnostic accuracy above 90\%. & $\square$ & $\square$ & $\square$ & $\square$ & $\square$ \\
\end{xltabular}}
\vspace{0.5em}\noindent\textit{Note.} SD = Strongly Disagree (1); D = Disagree (2); N = Neutral (3); A = Agree (4); SA = Strongly Agree (5). The 90\% accuracy threshold in B15 reflects the minimum acceptable performance identified in the RGAIG framework validation (Chapter~4).

%===============================================================================
%% SECTION C: ORGANIZATIONAL READINESS — TOE (12 items)
%===============================================================================
\subsection{Section C: Organizational Readiness (TOE Framework)}
\label{subsec:survey_toe}

\noindent\textit{Instructions: The following statements assess your organization's readiness to adopt AI-driven diagnostic tools. Please respond based on your current institutional context. If you are a patient or individual user, please answer based on the healthcare institution where you receive care, or select ``Not Applicable'' where appropriate.}

\needspace{24\baselineskip}
\noindent\textbf{C1--C4: Technology Context}

\noindent\textit{Theoretical basis: TOE Technology Context---the internal and external technologies relevant to the organization, including equipment and processes \cite{tornatzky1990processes}.}

\noindent Table~\ref{tab:survey_section_c_tech} reports the four Technology Context items evaluating institutional computing infrastructure, data maturity, system interoperability, and scalability.
{\tablestripes\tablefontsize
\needspace{20\baselineskip}
\begin{xltabular}{\textwidth}{@{} p{1.0cm} p{7.5cm} p{1.0cm} p{1.0cm} p{1.0cm} p{1.0cm} p{1.0cm} @{}}
\caption[Section C1--C4]{Section C1--C4: Technology Context Items}\label{tab:survey_section_c_tech} \\
\toprule
\thead{ID} & \thead{Statement} & \thead{SD} & \thead{D} & \thead{N} & \thead{A} & \thead{SA} \\
\midrule
\endfirsthead
\endhead
\bottomrule
\endlastfoot
C1 & My organization has the computing infrastructure (GPU servers, cloud resources, network bandwidth) necessary to conceptually evaluate AI-based diagnostic systems at clinical scale. & $\square$ & $\square$ & $\square$ & $\square$ & $\square$ \\
\addlinespace
C2 & My organization has sufficient high-quality, digitized clinical data (EEG recordings, imaging data, structured clinical notes) to support AI model training and validation. & $\square$ & $\square$ & $\square$ & $\square$ & $\square$ \\
\addlinespace
C3 & The existing clinical information systems at my organization (EHR, PACS, laboratory systems) are compatible with AI integration through standard APIs and data interoperability protocols (HL7 FHIR, DICOM). & $\square$ & $\square$ & $\square$ & $\square$ & $\square$ \\
\addlinespace
C4 & My organization has the ability to scale conceptual AI governance evaluations from a pilot (single department) to enterprise-wide adoption across multiple clinical departments. & $\square$ & $\square$ & $\square$ & $\square$ & $\square$ \\
\end{xltabular}}
\vspace{0.5em}\noindent\textit{Note.} SD = Strongly Disagree (1); D = Disagree (2); N = Neutral (3); A = Agree (4); SA = Strongly Agree (5). GPU = Graphics Processing Unit; EHR = Electronic Health Record; PACS = Picture Archiving and Communication System; HL7 FHIR = Health Level Seven Fast Healthcare Interoperability Resources; DICOM = Digital Imaging and Communications in Medicine; API = Application Programming Interface.

\needspace{24\baselineskip}
\noindent\textbf{C5--C8: Organization Context}
\paragraph{Section C5--C8.}
\noindent Table~\ref{tab:survey_section_c_org} below presents Section C5--C8, laying out each row in structured form to help examiners trace the source, applicability, and downstream use at a glance. % [unified-fix inserted]


\noindent\textit{Theoretical basis: TOE Organization Context---internal characteristics including management structure, communication processes, human resources, and slack \cite{tornatzky1990processes}.}

\noindent Table~\ref{tab:survey_section_c_org} presents the four Organization Context items measuring leadership support, budget allocation, training capacity, and change management readiness.
{\tablestripes\tablefontsize

\noindent Table~\ref{tab:survey_section_c_org} below presents Section C5--C8, laying out each row in structured form to help examiners trace the source, applicability, and downstream use at a glance. % [auto-intro-strict inserted]
\paragraph{Section C5--C8.} % [auto-heading inserted]
\noindent Table~\ref{tab:survey_section_c_org} below presents Section C5--C8, laying out each row in structured form to help examiners trace the source, applicability, and downstream use at a glance. % [auto-intro-after-heading]



\needspace{20\baselineskip}
\begin{xltabular}{\textwidth}{@{} p{1.0cm} p{7.5cm} p{1.0cm} p{1.0cm} p{1.0cm} p{1.0cm} p{1.0cm} @{}}
\caption[Section C5--C8]{Section C5--C8: Organization Context Items}\label{tab:survey_section_c_org} \\
\toprule
\thead{ID} & \thead{Statement} & \thead{SD} & \thead{D} & \thead{N} & \thead{A} & \thead{SA} \\
\midrule
\endfirsthead
\endhead
\bottomrule
\endlastfoot
C5 & Senior leadership at my organization (C-suite, department heads) actively champions and supports the adoption of AI-based clinical tools. & $\square$ & $\square$ & $\square$ & $\square$ & $\square$ \\
\addlinespace
C6 & My organization has allocated or is willing to allocate sufficient budget for AI diagnostic tool procurement, implementation, and ongoing maintenance. & $\square$ & $\square$ & $\square$ & $\square$ & $\square$ \\
\addlinespace
C7 & My organization has the capacity to provide adequate training programs for clinicians on how to use and interpret AI-assisted diagnostic tools. & $\square$ & $\square$ & $\square$ & $\square$ & $\square$ \\
\addlinespace
C8 & My organization has effective change management processes to facilitate the transition from traditional diagnostic workflows to AI-augmented workflows. & $\square$ & $\square$ & $\square$ & $\square$ & $\square$ \\
\end{xltabular}}
\vspace{0.5em}\noindent\textit{Note.} SD = Strongly Disagree (1); D = Disagree (2); N = Neutral (3); A = Agree (4); SA = Strongly Agree (5).

\needspace{24\baselineskip}
\noindent\textbf{C9--C12: Environment Context}
\paragraph{Section C9--C12.}
\noindent Table~\ref{tab:survey_section_c_env} below presents Section C9--C12, laying out each row in structured form to help examiners trace the source, applicability, and downstream use at a glance. % [unified-fix inserted]


\noindent\textit{Theoretical basis: TOE Environment Context---external factors including industry structure, regulatory environment, and competitive pressures \cite{tornatzky1990processes}.}

\noindent Table~\ref{tab:survey_section_c_env} documents the four Environment Context items assessing regulatory pressure, competitive dynamics, patient expectations, and industry standards influencing AI adoption.
{\tablestripes\tablefontsize

\noindent Table~\ref{tab:survey_section_c_env} below presents Section C9--C12, laying out each row in structured form to help examiners trace the source, applicability, and downstream use at a glance. % [auto-intro-strict inserted]
\paragraph{Section C9--C12.} % [auto-heading inserted]
\noindent Table~\ref{tab:survey_section_c_env} below presents Section C9--C12, laying out each row in structured form to help examiners trace the source, applicability, and downstream use at a glance. % [auto-intro-after-heading]



\needspace{20\baselineskip}
\begin{xltabular}{\textwidth}{@{} p{1.0cm} p{7.5cm} p{1.0cm} p{1.0cm} p{1.0cm} p{1.0cm} p{1.0cm} @{}}
\caption[Section C9--C12]{Section C9--C12: Environment Context Items}\label{tab:survey_section_c_env} \\
\toprule
\thead{ID} & \thead{Statement} & \thead{SD} & \thead{D} & \thead{N} & \thead{A} & \thead{SA} \\
\midrule
\endfirsthead
\endhead
\bottomrule
\endlastfoot
C9 & Evolving regulatory requirements (FDA SaMD guidelines, EU MDR, EU AI Act) are creating pressure for my organization to adopt validated AI diagnostic tools rather than informal or ad hoc approaches. & $\square$ & $\square$ & $\square$ & $\square$ & $\square$ \\
\addlinespace
C10 & Competing healthcare institutions in my region are already adopting or piloting AI-based diagnostic tools, creating competitive pressure for my organization to do the same. & $\square$ & $\square$ & $\square$ & $\square$ & $\square$ \\
\addlinespace
C11 & Patients increasingly expect or request that their healthcare providers use the latest AI-assisted diagnostic technologies. & $\square$ & $\square$ & $\square$ & $\square$ & $\square$ \\
\addlinespace
C12 & Established industry standards and best practice guidelines (e.g., NIST AI RMF, WHO AI guidance, AAN clinical practice guidelines) support the adoption of AI in neurological diagnosis at my institution. & $\square$ & $\square$ & $\square$ & $\square$ & $\square$ \\
\end{xltabular}}
\vspace{0.5em}\noindent\textit{Note.} SD = Strongly Disagree (1); D = Disagree (2); N = Neutral (3); A = Agree (4); SA = Strongly Agree (5). FDA = US Food and Drug Administration; SaMD = Software as a Medical Device; EU MDR = European Union Medical Device Regulation; EU AI Act = European Union Artificial Intelligence Act; NIST AI RMF = National Institute of Standards and Technology AI Risk Management Framework; WHO = World Health Organization; AAN = American Academy of Neurology.

%===============================================================================
%% SECTION D: TRUST AND GOVERNANCE (12 items)
%===============================================================================
\subsection{Section D: Trust and Governance}
\label{subsec:survey_trust}

\noindent\textit{Instructions: The following statements assess your level of trust in AI diagnostic systems and your confidence in governance mechanisms that oversee their deployment. Please indicate your level of agreement with each statement. These items are relevant to all stakeholder groups.}

\needspace{24\baselineskip}
\noindent\textbf{D1--D3: Algorithmic Trust}

\noindent\textit{Theoretical basis: Trust in automation \cite{lee2004trust} and algorithmic trust \cite{shin2021xai}---confidence in the system's competence, benevolence, and integrity.}

\noindent Table~\ref{tab:survey_section_d_algo} presents the three Algorithmic Trust items measuring confidence in AI system reliability, transparency, and traceability.
{\tablestripes\tablefontsize
\needspace{20\baselineskip}
\begin{xltabular}{\textwidth}{@{} p{1.0cm} p{7.5cm} p{1.0cm} p{1.0cm} p{1.0cm} p{1.0cm} p{1.0cm} @{}}
\caption[Section D1--D3]{Section D1--D3: Algorithmic Trust Items}\label{tab:survey_section_d_algo} \\
\toprule
\thead{ID} & \thead{Statement} & \thead{SD} & \thead{D} & \thead{N} & \thead{A} & \thead{SA} \\
\midrule
\endfirsthead
\endhead
\bottomrule
\endlastfoot
D1 & I have confidence that an AI system achieving $\geq$90\% accuracy on validated datasets would provide reliable diagnostic support in real-world clinical settings. & $\square$ & $\square$ & $\square$ & $\square$ & $\square$ \\
\addlinespace
D2 & I trust AI diagnostic systems more when they provide transparent explanations of how they arrived at each diagnostic recommendation (e.g., feature importance via SHAP values, attention maps). & $\square$ & $\square$ & $\square$ & $\square$ & $\square$ \\
\addlinespace
D3 & The ability to trace an AI diagnosis back to specific input features and supporting literature (explainability) is essential for me to trust the system's recommendations. & $\square$ & $\square$ & $\square$ & $\square$ & $\square$ \\
\end{xltabular}}
\vspace{0.5em}\noindent\textit{Note.} SD = Strongly Disagree (1); D = Disagree (2); N = Neutral (3); A = Agree (4); SA = Strongly Agree (5). SHAP = SHapley Additive exPlanations.

\needspace{24\baselineskip}
\noindent\textbf{D4--D6: Data Governance}
\paragraph{Section D4--D6.}
\noindent Table~\ref{tab:survey_section_d_data} below presents Section D4--D6, laying out each row in structured form to help examiners trace the source, applicability, and downstream use at a glance. % [unified-fix inserted]


\noindent\textit{Theoretical basis: Institutional trust theory (Zucker, 1986) and information privacy concerns (Smith et al., 1996)---confidence in organizational data handling practices.}

\noindent Table~\ref{tab:survey_section_d_data} reports the three Data Governance items assessing confidence in patient data privacy, consent management, and de-identification procedures.
{\tablestripes\tablefontsize

\noindent Table~\ref{tab:survey_section_d_data} below presents Section D4--D6, laying out each row in structured form to help examiners trace the source, applicability, and downstream use at a glance. % [auto-intro-strict inserted]
\paragraph{Section D4--D6.} % [auto-heading inserted]
\noindent Table~\ref{tab:survey_section_d_data} below presents Section D4--D6, laying out each row in structured form to help examiners trace the source, applicability, and downstream use at a glance. % [auto-intro-after-heading]



\needspace{20\baselineskip}
\begin{xltabular}{\textwidth}{@{} p{1.0cm} p{7.5cm} p{1.0cm} p{1.0cm} p{1.0cm} p{1.0cm} p{1.0cm} @{}}
\caption[Section D4--D6]{Section D4--D6: Data Governance Items}\label{tab:survey_section_d_data} \\
\toprule
\thead{ID} & \thead{Statement} & \thead{SD} & \thead{D} & \thead{N} & \thead{A} & \thead{SA} \\
\midrule
\endfirsthead
\endhead
\bottomrule
\endlastfoot
D4 & I am confident that AI diagnostic systems deployed at my institution would adequately protect patient data privacy in accordance with HIPAA, GDPR, and applicable local regulations. & $\square$ & $\square$ & $\square$ & $\square$ & $\square$ \\
\addlinespace
D5 & I believe that current consent management practices (informed consent for AI-assisted diagnosis, opt-in/opt-out mechanisms) are sufficient to protect patient autonomy in AI-driven care. & $\square$ & $\square$ & $\square$ & $\square$ & $\square$ \\
\addlinespace
D6 & I am confident that the de-identification and anonymization procedures applied to clinical datasets used for AI training are robust enough to prevent re-identification of individual patients. & $\square$ & $\square$ & $\square$ & $\square$ & $\square$ \\
\end{xltabular}}
\vspace{0.5em}\noindent\textit{Note.} SD = Strongly Disagree (1); D = Disagree (2); N = Neutral (3); A = Agree (4); SA = Strongly Agree (5). HIPAA = Health Insurance Portability and Accountability Act; GDPR = General Data Protection Regulation.

\needspace{24\baselineskip}
\noindent\textbf{D7--D9: Ethical Considerations}
\paragraph{Section D7--D9.}
\noindent Table~\ref{tab:survey_section_d_ethics} below presents Section D7--D9, laying out each row in structured form to help examiners trace the source, applicability, and downstream use at a glance. % [unified-fix inserted]


\noindent\textit{Theoretical basis: Responsible AI principles \cite{floridi2018ai4people}---fairness, accountability, transparency, and human oversight in automated decision-making.}

\noindent Table~\ref{tab:survey_section_d_ethics} documents the three Ethical Considerations items addressing demographic bias testing, human decision authority, and patient notification rights.
{\tablestripes\tablefontsize

\noindent Table~\ref{tab:survey_section_d_ethics} below presents Section D7--D9, laying out each row in structured form to help examiners trace the source, applicability, and downstream use at a glance. % [auto-intro-strict inserted]
\paragraph{Section D7--D9.} % [auto-heading inserted]
\noindent Table~\ref{tab:survey_section_d_ethics} below presents Section D7--D9, laying out each row in structured form to help examiners trace the source, applicability, and downstream use at a glance. % [auto-intro-after-heading]



\needspace{20\baselineskip}
\begin{xltabular}{\textwidth}{@{} p{1.0cm} p{7.5cm} p{1.0cm} p{1.0cm} p{1.0cm} p{1.0cm} p{1.0cm} @{}}
\caption[Section D7--D9]{Section D7--D9: Ethical Considerations Items}\label{tab:survey_section_d_ethics} \\
\toprule
\thead{ID} & \thead{Statement} & \thead{SD} & \thead{D} & \thead{N} & \thead{A} & \thead{SA} \\
\midrule
\endfirsthead
\endhead
\bottomrule
\endlastfoot
D7 & I am confident that the AI diagnostic system has been tested for demographic bias (across age, gender, ethnicity) and performs equitably across diverse patient populations. & $\square$ & $\square$ & $\square$ & $\square$ & $\square$ \\
\addlinespace
D8 & It is essential that a qualified human clinician retains final decision-making authority over all AI-generated diagnoses, even when the system achieves high accuracy. & $\square$ & $\square$ & $\square$ & $\square$ & $\square$ \\
\addlinespace
D9 & Patients should be fully informed when AI tools are used in their diagnostic process and should have the right to request a purely human-performed diagnosis. & $\square$ & $\square$ & $\square$ & $\square$ & $\square$ \\
\end{xltabular}}
\vspace{0.5em}\noindent\textit{Note.} SD = Strongly Disagree (1); D = Disagree (2); N = Neutral (3); A = Agree (4); SA = Strongly Agree (5).

\needspace{24\baselineskip}
\noindent\textbf{D10--D12: Regulatory Compliance}
\paragraph{Section D10--D12.}
\noindent Table~\ref{tab:survey_section_d_reg} below presents Section D10--D12, laying out each row in structured form to help examiners trace the source, applicability, and downstream use at a glance. % [unified-fix inserted]


\noindent\textit{Theoretical basis: Institutional theory \cite{dimaggio1983iron} and regulatory compliance frameworks---organizational conformity to external regulatory mandates governing AI in healthcare.}

\noindent Table~\ref{tab:survey_section_d_reg} presents the three Regulatory Compliance items measuring alignment with FDA SaMD requirements, HIPAA policies, and EU AI Act obligations.
{\tablestripes\tablefontsize

\noindent Table~\ref{tab:survey_section_d_reg} below presents Section D10--D12, laying out each row in structured form to help examiners trace the source, applicability, and downstream use at a glance. % [auto-intro-strict inserted]
\paragraph{Section D10--D12.} % [auto-heading inserted]
\noindent Table~\ref{tab:survey_section_d_reg} below presents Section D10--D12, laying out each row in structured form to help examiners trace the source, applicability, and downstream use at a glance. % [auto-intro-after-heading]



\needspace{20\baselineskip}
\begin{xltabular}{\textwidth}{@{} p{1.0cm} p{7.5cm} p{1.0cm} p{1.0cm} p{1.0cm} p{1.0cm} p{1.0cm} @{}}
\caption[Section D10--D12]{Section D10--D12: Regulatory Compliance Items}\label{tab:survey_section_d_reg} \\
\toprule
\thead{ID} & \thead{Statement} & \thead{SD} & \thead{D} & \thead{N} & \thead{A} & \thead{SA} \\
\midrule
\endfirsthead
\endhead
\bottomrule
\endlastfoot
D10 & I believe the RGAIG framework's quality assurance processes (525-module taxonomy with 97.4--98.4\% pass rates) are sufficient to meet FDA Software as a Medical Device (SaMD) regulatory requirements. & $\square$ & $\square$ & $\square$ & $\square$ & $\square$ \\
\addlinespace
D11 & My organization has the policies and procedures in place to ensure HIPAA compliance when deploying AI diagnostic tools that process protected health information (PHI). & $\square$ & $\square$ & $\square$ & $\square$ & $\square$ \\
\addlinespace
D12 & The RGAIG framework's governance structure (10 cross-cutting AI governance pillars, responsible AI protocols) aligns with the requirements of the EU AI Act for high-risk AI systems in healthcare. & $\square$ & $\square$ & $\square$ & $\square$ & $\square$ \\
\end{xltabular}}
\vspace{0.5em}{\small\noindent\textit{Note.} SD = Strongly Disagree (1); D = Disagree (2); N = Neutral (3); A = Agree (4); SA = Strongly Agree (5). FDA = US Food and Drug Administration; SaMD = Software as a Medical Device; PHI = Protected Health Information; HIPAA = Health Insurance Portability and Accountability Act; EU AI Act = European Union Artificial Intelligence Act.}

%===============================================================================
%% SECTION E: CLINICAL UTILITY ASSESSMENT — DSR (10 items)
%===============================================================================
\subsection{Section E}
\label{subsec:survey_clinical}

\noindent\textit{Instructions: The following statements assess the clinical utility of the RGAIG framework based on Design Science Research (DSR) evaluation criteria: utility, quality, and efficacy. Clinicians should respond based on their direct diagnostic experience; non-clinicians should respond based on their understanding of clinical needs within their professional role.}

\needspace{24\baselineskip}
\noindent\textbf{E1--E3: Diagnostic Performance (DSR Efficacy)}

\noindent\textit{Theoretical basis: DSR artifact evaluation---efficacy criterion \cite{hevner2004design}. Does the artifact achieve the intended outcomes in the target environment?}

\noindent Table~\ref{tab:survey_section_e_diag} reports the three Diagnostic Performance items evaluating whether the framework's reported accuracy, sensitivity, and specificity meet clinical adoption thresholds.
{\tablestripes\tablefontsize
\needspace{20\baselineskip}
\begin{xltabular}{\textwidth}{@{} p{1.0cm} p{7.5cm} p{1.0cm} p{1.0cm} p{1.0cm} p{1.0cm} p{1.0cm} @{}}
\caption[Section E1--E3]{Section E1--E3: Diagnostic Performance Items}\label{tab:survey_section_e_diag} \\
\toprule
\thead{ID} & \thead{Statement} & \thead{SD} & \thead{D} & \thead{N} & \thead{A} & \thead{SA} \\
\midrule
\endfirsthead
\endhead
\bottomrule
\endlastfoot
E1 & The reported ASD-classification accuracy of 97.67\% meets the minimum performance threshold I would require before considering clinical adoption. & $\square$ & $\square$ & $\square$ & $\square$ & $\square$ \\
\addlinespace
E2 & The system's sensitivity (ability to correctly identify patients with the disease) is adequate for use as a clinical screening tool, given that missed diagnoses carry significant patient harm. & $\square$ & $\square$ & $\square$ & $\square$ & $\square$ \\
\addlinespace
E3 & The system's specificity (ability to correctly identify patients without the disease) is adequate to avoid excessive false-positive referrals, which could burden specialist clinics and cause patient anxiety. & $\square$ & $\square$ & $\square$ & $\square$ & $\square$ \\
\end{xltabular}}
\vspace{0.5em}\noindent\textit{Note.} SD = Strongly Disagree (1); D = Disagree (2); N = Neutral (3); A = Agree (4); SA = Strongly Agree (5). Accuracy figures are from the RGAIG framework validation (Chapter~4, Table~4.X). ASD = Autism Spectrum Disorder.

\needspace{24\baselineskip}
\noindent\textbf{E4--E6: Workflow Integration (DSR Utility)}
\paragraph{Section E4--E6.}
\noindent Table~\ref{tab:survey_section_e_wf} below presents Section E4--E6, laying out each row in structured form to help examiners trace the source, applicability, and downstream use at a glance. % [unified-fix inserted]


\noindent\textit{Theoretical basis: DSR artifact evaluation---utility criterion (Hevner et al., 2004; March \& Smith, 1995). Does the artifact provide value in the organizational context?}

\noindent Table~\ref{tab:survey_section_e_wf} presents the three Workflow Integration items assessing report quality, turnaround time, and clinical actionability of AI-generated diagnostic outputs.
{\tablestripes\tablefontsize

\noindent Table~\ref{tab:survey_section_e_wf} below presents Section E4--E6, laying out each row in structured form to help examiners trace the source, applicability, and downstream use at a glance. % [auto-intro-strict inserted]
\paragraph{Section E4--E6.} % [auto-heading inserted]
\noindent Table~\ref{tab:survey_section_e_wf} below presents Section E4--E6, laying out each row in structured form to help examiners trace the source, applicability, and downstream use at a glance. % [auto-intro-after-heading]



\needspace{20\baselineskip}
\begin{xltabular}{\textwidth}{@{} p{1.0cm} p{7.5cm} p{1.0cm} p{1.0cm} p{1.0cm} p{1.0cm} p{1.0cm} @{}}
\caption[Section E4--E6]{Section E4--E6: Workflow Integration Items}\label{tab:survey_section_e_wf} \\
\toprule
\thead{ID} & \thead{Statement} & \thead{SD} & \thead{D} & \thead{N} & \thead{A} & \thead{SA} \\
\midrule
\endfirsthead
\endhead
\bottomrule
\endlastfoot
E4 & The AI-generated diagnostic reports (including classification results, feature importance, and literature-grounded explanations) would be of sufficient quality and clarity to support my clinical documentation requirements. & $\square$ & $\square$ & $\square$ & $\square$ & $\square$ \\
\addlinespace
E5 & The system's expected turnaround time (under 2 minutes for a complete multi-modal analysis) would be fast enough to integrate into my standard clinical appointment workflow. & $\square$ & $\square$ & $\square$ & $\square$ & $\square$ \\
\addlinespace
E6 & The diagnostic outputs of the AI system would be clinically actionable---that is, they would lead directly to specific follow-up actions (referral, treatment initiation, further testing) without requiring additional interpretation. & $\square$ & $\square$ & $\square$ & $\square$ & $\square$ \\
\end{xltabular}}
\vspace{0.5em}\noindent\textit{Note.} SD = Strongly Disagree (1); D = Disagree (2); N = Neutral (3); A = Agree (4); SA = Strongly Agree (5).

\needspace{24\baselineskip}
\noindent\textbf{E7--E10: ASD-Focused Value (DSR Quality)}
\paragraph{Section E7--E10.}
\noindent Table~\ref{tab:survey_section_e_multi} below presents Section E7--E10, laying out each row in structured form to help examiners trace the source, applicability, and downstream use at a glance. % [unified-fix inserted]


\noindent\textit{Theoretical basis: DSR artifact evaluation---quality criterion \cite{hevner2004design}. RBV value attribute---does the ASD-focused capability represent a valuable, rare, and non-substitutable resource \cite{barney1991firm}?}

\noindent Table~\ref{tab:survey_section_e_multi} documents the four ASD-Focused Value items measuring ASD-focused comorbidity detection, rare presentation identification, second-opinion utility, and educational value.
{\tablestripes\tablefontsize

\noindent Table~\ref{tab:survey_section_e_multi} below presents Section E7--E10, laying out each row in structured form to help examiners trace the source, applicability, and downstream use at a glance. % [auto-intro-strict inserted]
\paragraph{Section E7--E10.} % [auto-heading inserted]
\noindent Table~\ref{tab:survey_section_e_multi} below presents Section E7--E10, laying out each row in structured form to help examiners trace the source, applicability, and downstream use at a glance. % [auto-intro-after-heading]



\needspace{20\baselineskip}
\begin{xltabular}{\textwidth}{@{} p{1.0cm} p{7.5cm} p{1.0cm} p{1.0cm} p{1.0cm} p{1.0cm} p{1.0cm} @{}}
\caption[Section E7--E10]{Section E7--E10: ASD-Focused Value Items}\label{tab:survey_section_e_multi} \\
\toprule
\thead{ID} & \thead{Statement} & \thead{SD} & \thead{D} & \thead{N} & \thead{A} & \thead{SA} \\
\midrule
\endfirsthead
\endhead
\bottomrule
\endlastfoot
E7 & A unified AI framework capable of screening across multiple neurological domains (rather than separate single-disease tools) would provide significant value for ASD-focused comorbidity detection. & $\square$ & $\square$ & $\square$ & $\square$ & $\square$ \\
\addlinespace
E8 & The AI system's ability to detect rare or atypical presentations of neurological conditions would be a valuable supplement to clinical expertise, particularly in settings with limited specialist access. & $\square$ & $\square$ & $\square$ & $\square$ & $\square$ \\
\addlinespace
E9 & I would value AI-generated diagnostic results as a ``second opinion'' to complement my own clinical judgment, particularly for complex or ambiguous cases. & $\square$ & $\square$ & $\square$ & $\square$ & $\square$ \\
\addlinespace
E10 & The RGAIG framework would serve as a valuable training and education tool for medical residents and junior clinicians learning neurological diagnosis. & $\square$ & $\square$ & $\square$ & $\square$ & $\square$ \\
\end{xltabular}}
\vspace{0.5em}\noindent\textit{Note.} SD = Strongly Disagree (1); D = Disagree (2); N = Neutral (3); A = Agree (4); SA = Strongly Agree (5). RGAIG = Responsible Generative AI Governance.

%===============================================================================
%% SECTION F: RAG SYSTEM EVALUATION (8 items)
%===============================================================================
\subsection{Section F: RAG System Evaluation}
\label{subsec:survey_rag}

\noindent\textit{Instructions: The RGAIG framework includes a Retrieval-Augmented Generation (RAG) engine that produces natural-language diagnostic explanations grounded in peer-reviewed medical literature. The RAG system retrieves relevant publications from a curated knowledge base and generates explanatory narratives that cite specific studies. Please indicate your level of agreement with the following statements about this component.}

\needspace{24\baselineskip}
\noindent\textbf{F1--F3: Report Quality (RBV Value and Rarity)}
\paragraph{Section F1--F3.}
\noindent Table~\ref{tab:survey_section_f_qual} below presents Section F1--F3, laying out each row in structured form to help examiners trace the source, applicability, and downstream use at a glance. % [unified-fix inserted]


\noindent\textit{Theoretical basis: RBV---the RAG-generated reports represent a valuable, rare, and difficult-to-imitate resource that provides competitive advantage \cite{barney1991firm}. DSR quality evaluation \cite{hevner2004design}.}

\noindent Table~\ref{tab:survey_section_f_qual} reports the three RAG Report Quality items evaluating factual faithfulness, citation verifiability, and clinical relevance of generated explanations.
{\tablestripes\tablefontsize

\noindent Table~\ref{tab:survey_section_f_qual} below presents Section F1--F3, laying out each row in structured form to help examiners trace the source, applicability, and downstream use at a glance. % [auto-intro-strict inserted]
\paragraph{Section F1--F3.} % [auto-heading inserted]
\noindent Table~\ref{tab:survey_section_f_qual} below presents Section F1--F3, laying out each row in structured form to help examiners trace the source, applicability, and downstream use at a glance. % [auto-intro-after-heading]



\needspace{20\baselineskip}
\begin{xltabular}{\textwidth}{@{} p{1.0cm} p{7.5cm} p{1.0cm} p{1.0cm} p{1.0cm} p{1.0cm} p{1.0cm} @{}}
\caption[Section F1--F3]{Section F1--F3: RAG Report Quality Items}\label{tab:survey_section_f_qual} \\
\toprule
\thead{ID} & \thead{Statement} & \thead{SD} & \thead{D} & \thead{N} & \thead{A} & \thead{SA} \\
\midrule
\endfirsthead
\endhead
\bottomrule
\endlastfoot
F1 & The RAG-generated diagnostic explanations would be factually faithful---that is, the content of the explanation would accurately reflect the underlying AI classification results and not introduce unsupported claims. & $\square$ & $\square$ & $\square$ & $\square$ & $\square$ \\
\addlinespace
F2 & The literature citations provided in RAG-generated reports (specific author names, publication years, journal references) would be verifiable and accurately attributed. & $\square$ & $\square$ & $\square$ & $\square$ & $\square$ \\
\addlinespace
F3 & The RAG-generated explanations would be clinically relevant---that is, they would reference literature and findings that are directly pertinent to the patient's specific diagnostic context. & $\square$ & $\square$ & $\square$ & $\square$ & $\square$ \\
\end{xltabular}}
\vspace{0.5em}\noindent\textit{Note.} SD = Strongly Disagree (1); D = Disagree (2); N = Neutral (3); A = Agree (4); SA = Strongly Agree (5). RAG = Retrieval-Augmented Generation.

\needspace{24\baselineskip}
\noindent\textbf{F4--F6: Explainability Value (RBV Inimitability)}
\paragraph{Section F4--F6.}
\noindent Table~\ref{tab:survey_section_f_xai} below presents Section F4--F6, laying out each row in structured form to help examiners trace the source, applicability, and downstream use at a glance. % [unified-fix inserted]


\noindent\textit{Theoretical basis: RBV inimitability---the combination of SHAP/LIME feature importance with RAG-generated literature grounding is a complex, path-dependent capability that is difficult for competitors to replicate \cite{barney1991firm}. Explainable AI (XAI) principles \cite{arrieta2020xai}.}

\noindent Table~\ref{tab:survey_section_f_xai} presents the three Explainability Value items measuring the trust impact of SHAP/LIME visualisations, literature grounding, and confidence calibration.
{\tablestripes\tablefontsize

\noindent Table~\ref{tab:survey_section_f_xai} below presents Section F4--F6, laying out each row in structured form to help examiners trace the source, applicability, and downstream use at a glance. % [auto-intro-strict inserted]
\paragraph{Section F4--F6.} % [auto-heading inserted]
\noindent Table~\ref{tab:survey_section_f_xai} below presents Section F4--F6, laying out each row in structured form to help examiners trace the source, applicability, and downstream use at a glance. % [auto-intro-after-heading]



\needspace{20\baselineskip}
\begin{xltabular}{\textwidth}{@{} p{1.0cm} p{7.5cm} p{1.0cm} p{1.0cm} p{1.0cm} p{1.0cm} p{1.0cm} @{}}
\caption[Section F4--F6]{Section F4--F6: Explainability Value Items}\label{tab:survey_section_f_xai} \\
\toprule
\thead{ID} & \thead{Statement} & \thead{SD} & \thead{D} & \thead{N} & \thead{A} & \thead{SA} \\
\midrule
\endfirsthead
\endhead
\bottomrule
\endlastfoot
F4 & Feature importance visualizations (SHAP values, LIME explanations) that show which EEG channels, frequency bands, or imaging features contributed most to the diagnosis would significantly increase my trust in the AI output. & $\square$ & $\square$ & $\square$ & $\square$ & $\square$ \\
\addlinespace
F5 & Grounding AI diagnostic explanations in specific peer-reviewed literature (rather than generic statistical summaries) would make the output more credible and useful for clinical communication with patients and colleagues. & $\square$ & $\square$ & $\square$ & $\square$ & $\square$ \\
\addlinespace
F6 & The system's reported confidence scores (probability estimates for each diagnostic category) would need to be well-calibrated---that is, a 90\% confidence score should correspond to approximately 90\% actual accuracy---for me to rely on them clinically. & $\square$ & $\square$ & $\square$ & $\square$ & $\square$ \\
\end{xltabular}}
\vspace{0.5em}\noindent\textit{Note.} SD = Strongly Disagree (1); D = Disagree (2); N = Neutral (3); A = Agree (4); SA = Strongly Agree (5). SHAP = SHapley Additive exPlanations; LIME = Local Interpretable Model-agnostic Explanations; XAI = Explainable Artificial Intelligence.

\needspace{24\baselineskip}
\noindent\textbf{F7--F8: Hallucination Concerns}
\paragraph{Section F7--F8.}
\noindent Table~\ref{tab:survey_section_f_halluc} below presents Section F7--F8, laying out each row in structured form to help examiners trace the source, applicability, and downstream use at a glance. % [unified-fix inserted]


\noindent\textit{Theoretical basis: Generative AI hallucination risk (Ji et al., 2023)---the tendency of large language models to generate plausible but factually incorrect content, which poses unique risks in clinical settings.}

\noindent Table~\ref{tab:survey_section_f_halluc} documents the two Hallucination Concern items assessing the extent to which fabricated citations and incorrect claims undermine trust and whether the RGAIG mitigation strategies are perceived as adequate.
{\tablestripes\tablefontsize

\noindent Table~\ref{tab:survey_section_f_halluc} below presents Section F7--F8, laying out each row in structured form to help examiners trace the source, applicability, and downstream use at a glance. % [auto-intro-strict inserted]
\paragraph{Section F7--F8.} % [auto-heading inserted]
\noindent Table~\ref{tab:survey_section_f_halluc} below presents Section F7--F8, laying out each row in structured form to help examiners trace the source, applicability, and downstream use at a glance. % [auto-intro-after-heading]



\needspace{20\baselineskip}
\begin{xltabular}{\textwidth}{@{} p{1.0cm} p{7.5cm} p{1.0cm} p{1.0cm} p{1.0cm} p{1.0cm} p{1.0cm} @{}}
\caption[Section F7--F8]{Section F7--F8: Hallucination Concern Items}\label{tab:survey_section_f_halluc} \\
\toprule
\thead{ID} & \thead{Statement} & \thead{SD} & \thead{D} & \thead{N} & \thead{A} & \thead{SA} \\
\midrule
\endfirsthead
\endhead
\bottomrule
\endlastfoot
F7 & The risk of AI-generated hallucinations (fabricated citations, incorrect clinical claims, spurious diagnostic rationales) is a significant barrier to my trust in RAG-generated diagnostic reports. & $\square$ & $\square$ & $\square$ & $\square$ & $\square$ \\
\addlinespace
F8 & I am confident that the RGAIG framework's hallucination mitigation strategies (retrieval grounding, citation verification, confidence thresholds, 13-phase GenAI quality assurance) are adequate to reduce hallucination risk to clinically acceptable levels. & $\square$ & $\square$ & $\square$ & $\square$ & $\square$ \\
\end{xltabular}}
\vspace{0.5em}\noindent\textit{Note.} SD = Strongly Disagree (1); D = Disagree (2); N = Neutral (3); A = Agree (4); SA = Strongly Agree (5). RAG = Retrieval-Augmented Generation. Item F7 is reverse-scored: higher agreement indicates greater concern (lower trust). The 13-phase GenAI quality assurance taxonomy (220 modules, 98.1\% pass rate) is detailed in Chapter~3.

%===============================================================================
%% SECTION G: IMPLEMENTATION BARRIERS (7 items, open-ended + ranked)
%===============================================================================
\subsection{Section G}
\label{subsec:survey_barriers}

\noindent\textit{Instructions: This section assesses the barriers, requirements, and expectations that would need to be addressed before you would support or participate in the future clinical-validation pathway of an AI-driven neurological diagnostic system. Please provide thoughtful responses; these qualitative data will be analyzed thematically to complement the quantitative findings.}

\noindent Table~\ref{tab:survey_section_g} presents the seven Implementation Barriers items combining ranked, categorical, and open-ended response formats to capture adoption obstacles, regulatory requirements, accuracy thresholds, and privacy concerns.
{\tablestripes\tablefontsize
\needspace{20\baselineskip}
\begin{xltabular}{\textwidth}{@{} p{1.0cm} p{8.2cm} p{5.2cm} @{}}
\caption[Section G]{Section G: Implementation Barriers and Adoption Requirements (G1--G7)}\label{tab:survey_section_g} \\
\toprule
\thead{ID} & \thead{Question} & \thead{Response Format} \\
\midrule
\endfirsthead
\endhead
\bottomrule
\endlastfoot
G1 & Please rank the top 3 barriers to adopting an AI-driven neurological diagnostic system at your institution from the following list: (a)~insufficient evidence of clinical efficacy, (b)~data privacy and security concerns, (c)~lack of regulatory approval (FDA/CE marking), (d)~high implementation cost, (e)~clinician resistance to change, (f)~lack of IT infrastructure, (g)~liability and malpractice uncertainty, (h)~workflow disruption, (i)~algorithmic bias concerns, (j)~lack of explainability. & Rank 1st, 2nd, 3rd: \par 1st: \rule{4cm}{0.4pt} \par 2nd: \rule{4cm}{0.4pt} \par 3rd: \rule{4cm}{0.4pt} \\
\addlinespace
G2 & What specific regulatory approvals or certifications would your institution require before deploying an AI diagnostic tool in clinical settings? (Select all that apply and/or specify others.) & $\square$~FDA 510(k) clearance \par $\square$~FDA De Novo classification \par $\square$~CE marking (EU MDR) \par $\square$~IRB approval for clinical use \par $\square$~Institutional future clinical-validation pathway study \par $\square$~Professional society endorsement (e.g., AAN) \par $\square$~Other: \rule{3cm}{0.4pt} \\
\addlinespace
G3 & What is the minimum diagnostic accuracy threshold (overall percentage) that an AI system must achieve before you would consider it suitable for clinical adoption as a screening or decision-support tool? & $\square$~$\geq$80\% \quad $\square$~$\geq$85\% \quad $\square$~$\geq$90\% \quad $\square$~$\geq$95\% \quad $\square$~$\geq$99\% \par Rationale (optional): \rule{4cm}{0.4pt} \\
\addlinespace
G4 & What is the maximum acceptable training duration for clinicians to become proficient in using the AI diagnostic system? & $\square$~Less than 1 day \par $\square$~1--3 days \par $\square$~1--2 weeks \par $\square$~2--4 weeks \par $\square$~More than 1 month \\
\addlinespace
G5 & Please describe your primary data privacy concern(s) regarding the use of AI systems that process sensitive neurological data (EEG, brain imaging, clinical records). & Open-ended: \par \rule{\linewidth}{0.4pt} \par \rule{\linewidth}{0.4pt} \par \rule{\linewidth}{0.4pt} \\
\addlinespace
G6 & Within what timeframe would you expect an AI diagnostic system to demonstrate a positive return on investment (ROI) for your institution? & $\square$~Within 6 months \par $\square$~6--12 months \par $\square$~1--2 years \par $\square$~2--3 years \par $\square$~More than 3 years \par $\square$~ROI is not a primary consideration \\
\addlinespace
G7 & Please provide any additional comments, concerns, or suggestions regarding the adoption of AI-driven diagnostic tools for neurological disease detection. & Open-ended: \par \rule{\linewidth}{0.4pt} \par \rule{\linewidth}{0.4pt} \par \rule{\linewidth}{0.4pt} \par \rule{\linewidth}{0.4pt} \\
\end{xltabular}}
\vspace{0.5em}\noindent\textit{Note.} Items G1, G3, G4, and G6 are categorical or ranked. Items G2 is multi-select. Items G5 and G7 are open-ended and will be analyzed using reflexive thematic analysis \cite{braun2006thematic}. FDA = US Food and Drug Administration; CE = Conformit\'{e} Europ\'{e}enne; EU MDR = European Union Medical Device Regulation; IRB = Institutional Review Board; AAN = American Academy of Neurology; illustrative financial scenario = return on investment.

%===============================================================================
%% SECTION H: WEARABLE INTEGRATION (8 items)
%===============================================================================
\subsection{Section H}
\label{subsec:survey_wearable}

\noindent\textit{Instructions: The RGAIG framework supports integration with consumer-grade wearable EEG devices (e.g., Emotiv Insight, Emotiv EPOC X) and multi-modal wearable sensor arrays (ECG, PPG, EDA, IMU, temperature) for continuous neurological monitoring outside clinical settings. Please indicate your level of agreement with the following statements about this wearable integration capability.}

\needspace{24\baselineskip}
\noindent\textbf{H1--H3: Consumer Wearable EEG Acceptance}
\paragraph{Section H1--H3.}
\noindent Table~\ref{tab:survey_section_h_eeg} below presents Section H1--H3, laying out each row in structured form to help examiners trace the source, applicability, and downstream use at a glance. % [unified-fix inserted]


\noindent\textit{Theoretical basis: TAM applied to wearable health technology (Kim \& Park, 2012)---perceived usefulness and ease of use of consumer wearable devices for clinical data collection.}

\noindent Table~\ref{tab:survey_section_h_eeg} reports the three Consumer Wearable EEG Acceptance items assessing perceptions of data quality from consumer-grade devices, longitudinal monitoring value, and remote care continuity.
{\tablestripes\tablefontsize

\noindent Table~\ref{tab:survey_section_h_eeg} below presents Section H1--H3, laying out each row in structured form to help examiners trace the source, applicability, and downstream use at a glance. % [auto-intro-strict inserted]
\paragraph{Section H1--H3.} % [auto-heading inserted]
\noindent Table~\ref{tab:survey_section_h_eeg} below presents Section H1--H3, laying out each row in structured form to help examiners trace the source, applicability, and downstream use at a glance. % [auto-intro-after-heading]



\needspace{20\baselineskip}
\begin{xltabular}{\textwidth}{@{} p{1.0cm} p{7.5cm} p{1.0cm} p{1.0cm} p{1.0cm} p{1.0cm} p{1.0cm} @{}}
\caption[Section H1--H3]{Section H1--H3: Consumer Wearable EEG Acceptance Items}\label{tab:survey_section_h_eeg} \\
\toprule
\thead{ID} & \thead{Statement} & \thead{SD} & \thead{D} & \thead{N} & \thead{A} & \thead{SA} \\
\midrule
\endfirsthead
\endhead
\bottomrule
\endlastfoot
H1 & Consumer-grade EEG devices (e.g., Emotiv Insight with 5 channels) can provide data of sufficient quality for AI-driven neurological screening, despite having fewer channels than clinical-grade systems (64--256 channels). & $\square$ & $\square$ & $\square$ & $\square$ & $\square$ \\
\addlinespace
H2 & Enabling patients to collect EEG data at home using consumer wearable devices would be valuable for longitudinal monitoring of neurological conditions between clinical visits. & $\square$ & $\square$ & $\square$ & $\square$ & $\square$ \\
\addlinespace
H3 & Remote monitoring through wearable devices, with AI-analyzed results transmitted to the clinical team in near-real-time, would improve continuity of care for patients with chronic neurological conditions. & $\square$ & $\square$ & $\square$ & $\square$ & $\square$ \\
\end{xltabular}}
\vspace{0.5em}\noindent\textit{Note.} SD = Strongly Disagree (1); D = Disagree (2); N = Neutral (3); A = Agree (4); SA = Strongly Agree (5). EEG = electroencephalography.

\needspace{24\baselineskip}
\noindent\textbf{H4--H6: Sensor Fusion Value}
\paragraph{Section H4--H6.}
\noindent Table~\ref{tab:survey_section_h_fusion} below presents Section H4--H6, laying out each row in structured form to help examiners trace the source, applicability, and downstream use at a glance. % [unified-fix inserted]


\noindent\textit{Theoretical basis: RBV non-substitutability---the combination of multi-modal sensor data (EEG + ECG + EDA + accelerometry) creates a complex, non-substitutable diagnostic resource \cite{barney1991firm}. DSR utility evaluation \cite{hevner2004design}.}

\noindent Table~\ref{tab:survey_section_h_fusion} presents the three Sensor Fusion Value items measuring the perceived diagnostic benefit of multi-modal sensor combination, continuous monitoring, and configurable alert thresholds.
{\tablestripes\tablefontsize

\noindent Table~\ref{tab:survey_section_h_fusion} below presents Section H4--H6, laying out each row in structured form to help examiners trace the source, applicability, and downstream use at a glance. % [auto-intro-strict inserted]
\paragraph{Section H4--H6.} % [auto-heading inserted]
\noindent Table~\ref{tab:survey_section_h_fusion} below presents Section H4--H6, laying out each row in structured form to help examiners trace the source, applicability, and downstream use at a glance. % [auto-intro-after-heading]



\needspace{20\baselineskip}
\begin{xltabular}{\textwidth}{@{} p{1.0cm} p{7.5cm} p{1.0cm} p{1.0cm} p{1.0cm} p{1.0cm} p{1.0cm} @{}}
\caption[Section H4--H6]{Section H4--H6: Sensor Fusion Value Items}\label{tab:survey_section_h_fusion} \\
\toprule
\thead{ID} & \thead{Statement} & \thead{SD} & \thead{D} & \thead{N} & \thead{A} & \thead{SA} \\
\midrule
\endfirsthead
\endhead
\bottomrule
\endlastfoot
H4 & Combining data from multiple sensor modalities (EEG, ECG, EDA, accelerometry, temperature) would provide richer diagnostic information than any single data source alone. & $\square$ & $\square$ & $\square$ & $\square$ & $\square$ \\
\addlinespace
H5 & Continuous wearable monitoring (24/7 data collection with AI analysis) would enable earlier detection of disease progression or treatment response than periodic clinic-based assessments. & $\square$ & $\square$ & $\square$ & $\square$ & $\square$ \\
\addlinespace
H6 & Configurable alert thresholds (e.g., automatic clinician notification when EEG patterns deviate beyond a defined baseline by more than 2 standard deviations) would be a valuable safety feature for remote patient monitoring. & $\square$ & $\square$ & $\square$ & $\square$ & $\square$ \\
\end{xltabular}}
\vspace{0.5em}\noindent\textit{Note.} SD = Strongly Disagree (1); D = Disagree (2); N = Neutral (3); A = Agree (4); SA = Strongly Agree (5). EEG = electroencephalography; ECG = electrocardiography; EDA = electrodermal activity; IMU = Inertial Measurement Unit.

\needspace{24\baselineskip}
\noindent\textbf{H7--H8: Patient Experience}
\paragraph{Section H7--H8.}
\noindent Table~\ref{tab:survey_section_h_patient} below presents Section H7--H8, laying out each row in structured form to help examiners trace the source, applicability, and downstream use at a glance. % [unified-fix inserted]


\noindent\textit{Theoretical basis: TAM Perceived Ease of Use extended to patient-facing technology \cite{holden2010tam}---patient comfort, compliance, and willingness to use wearable health monitoring devices.}

\noindent Table~\ref{tab:survey_section_h_patient} documents the two Patient Experience items assessing wearable comfort for extended use and patient willingness to comply with prescribed monitoring protocols.
{\tablestripes\tablefontsize

\noindent Table~\ref{tab:survey_section_h_patient} below presents Section H7--H8, laying out each row in structured form to help examiners trace the source, applicability, and downstream use at a glance. % [auto-intro-strict inserted]
\paragraph{Section H7--H8.} % [auto-heading inserted]
\noindent Table~\ref{tab:survey_section_h_patient} below presents Section H7--H8, laying out each row in structured form to help examiners trace the source, applicability, and downstream use at a glance. % [auto-intro-after-heading]



\needspace{20\baselineskip}
\begin{xltabular}{\textwidth}{@{} p{1.0cm} p{7.5cm} p{1.0cm} p{1.0cm} p{1.0cm} p{1.0cm} p{1.0cm} @{}}
\caption[Section H7--H8]{Section H7--H8: Patient Experience Items}\label{tab:survey_section_h_patient} \\
\toprule
\thead{ID} & \thead{Statement} & \thead{SD} & \thead{D} & \thead{N} & \thead{A} & \thead{SA} \\
\midrule
\endfirsthead
\endhead
\bottomrule
\endlastfoot
H7 & Consumer-grade wearable EEG devices are comfortable enough for patients (including elderly patients and children with neurological conditions) to wear for extended periods (4 or more hours per day). & $\square$ & $\square$ & $\square$ & $\square$ & $\square$ \\
\addlinespace
H8 & Patients (or their caregivers) would be willing to comply with a prescribed wearable monitoring protocol (e.g., daily 2-hour EEG recording sessions over a 4-week period) if they understood it would improve their diagnostic accuracy and treatment monitoring. & $\square$ & $\square$ & $\square$ & $\square$ & $\square$ \\
\end{xltabular}}
\vspace{0.5em}\noindent\textit{Note.} SD = Strongly Disagree (1); D = Disagree (2); N = Neutral (3); A = Agree (4); SA = Strongly Agree (5).

\subsubsection{B2C-Primary Adoption Addendum}
\label{subsec:survey_b2c_primary_addendum}

\noindent Because this dissertation positions B2C caregiver/patient use as the primary value pathway, the survey requires a short B2C addendum that directly tests adoption, trust, social setting, affordability, country context, and conceptual escalation reference (future scope). These items should be administered to patients, caregivers, and family respondents, and may be analysed descriptively if the B2C sample is too small for full SEM.

\noindent Table~\ref{tab:survey_b2c_primary_addendum} summarises b2c-primary adoption addendum for appendix review.

{\tablestripes\tablefontsize
\needspace{20\baselineskip}
\begin{xltabular}{\textwidth}{@{} p{1.0cm} p{9.5cm} p{3.5cm} @{}}
\caption[B2C-primary adoption addendum]{B2C-Primary Adoption Addendum. These items convert the consumer value proposition into testable survey evidence for caregiver/patient adoption.}\label{tab:survey_b2c_primary_addendum} \\
\toprule
\thead{Item} & \thead{Question} & \thead{Response} \\
\midrule
\endfirsthead
\endhead
\bottomrule
\endlastfoot
B2C1 & I would use this device at home if the result was clearly marked as screening support, not final diagnosis. & Likert 1--5 \\
B2C2 & I would trust the result more if a clinician could review the report before any major decision. & Likert 1--5 \\
B2C3 & Continuous monitoring would help me understand behavioural or neurological changes between clinic visits. & Likert 1--5 \\
B2C4 & I would allow conceptual escalation reference (future scope) to a hospital or clinician for critical alerts if the consent rules were clear. & Likert 1--5 \\
B2C5 & My main concern would be privacy of child/patient brain and behaviour data. & Likert 1--5 \\
B2C6 & The device must be comfortable and discreet enough for home, school, or public settings. & Likert 1--5 \\
B2C7 & I would pay for this service only if it reduced waiting time, travel, or repeated clinic visits. & Likert 1--5 \\
B2C8 & Country/location context: Canada, India, United States, or other. & Single select \\
B2C9 & User pathway: child caregiver, teen caregiver, adult self-user, adult caregiver, or clinician-assisted family. & Single select \\
B2C10 & What would stop you from using this device? & Open text \\
\end{xltabular}}
\vspace{0.5em}\noindent\textit{Note.} The addendum should be reported separately from clinician adoption results. Its purpose is to test consumer trust, affordability, comfort, privacy, and escalation acceptance before any public B2C launch.

%===============================================================================
%% CONSTRUCT-QUESTION-HYPOTHESIS-ANALYSIS MAPPING TABLE
%===============================================================================
\subsection{Construct--Question--Hypothesis--Analysis Mapping}
\label{subsec:survey_mapping}

Table~\ref{tab:survey_mapping} provides a complete traceability matrix linking each survey item to its parent construct, the theoretical framework, the research hypothesis it informs, and the planned statistical analysis method.
{\tablestripes\tablefontsize

\noindent Table~\ref{tab:survey_mapping} presents the survey item traceability: question to construct to.

\needspace{20\baselineskip}
\begin{xltabular}{\textwidth}{@{} p{1.0cm} L L L L @{}}
\caption[Survey Item Traceability: Question to Construct to]{Survey Item Traceability: Question to Construct to Hypothesis to Analysis Method}\label{tab:survey_mapping} \\
\toprule
\thead{Items} & \thead{Construct} & \thead{Theory} & \thead{Hypothesis} & \thead{Analysis Method} \\
\midrule
\endfirsthead
\endhead
\bottomrule
\endlastfoot
B1--B6 & Perceived Usefulness (PU) & TAM & H1: PU $\rightarrow$ BI & PLS-SEM path coefficient ($\beta$, $p < .05$) \\
\addlinespace
B7--B12 & Perceived Ease of Use (PEOU) & TAM & H2: PEOU $\rightarrow$ PU; H3: PEOU $\rightarrow$ BI & PLS-SEM path coefficient; mediation \\
\addlinespace
B13--B15 & Behavioural Intention (BI) & TAM & DV for H1--H3 & $R^2$, predictive relevance ($Q^2$) \\
\addlinespace
C1--C4 & Technology Context (TC) & TOE & H4: TC $\rightarrow$ Adoption readiness & PLS-SEM; MGA by institution type \\
\addlinespace
C5--C8 & Organization Context (OC) & TOE & H5: OC moderates PU $\rightarrow$ BI & Moderation (interaction term) \\
\addlinespace
C9--C12 & Environment Context (EC) & TOE & H6: EC $\rightarrow$ Adoption urgency & PLS-SEM; control for region (A8) \\
\addlinespace
D1--D3 & Algorithmic Trust (AT) & Trust Theory & H7: AT mediates PU $\rightarrow$ BI & Mediation (bootstrapped indirect effect) \\
\addlinespace
D4--D6 & Data Governance (DG) & Institutional Trust & H8: DG $\rightarrow$ AT & PLS-SEM path coefficient \\
\addlinespace
D7--D9 & Ethical Considerations (ETH) & Responsible AI & H9: ETH moderates AT $\rightarrow$ BI & Moderation analysis \\
\addlinespace
D10--D12 & Regulatory Compliance (RC) & Institutional Theory & H10: RC $\rightarrow$ OC & PLS-SEM path coefficient \\
\addlinespace
E1--E3 & Diagnostic Performance (DP) & DSR Efficacy & H11: DP $\rightarrow$ PU & PLS-SEM; IPMA \\
\addlinespace
E4--E6 & Workflow Integration (WI) & DSR Utility & H12: WI $\rightarrow$ PEOU & PLS-SEM path coefficient \\
\addlinespace
E7--E10 & ASD-Focused Value (MDV) & DSR Quality + RBV & H13: MDV $\rightarrow$ PU; H14: VRIN & PLS-SEM; RBV evaluation (descriptive) \\
\addlinespace
F1--F3 & RAG Report Quality (RQ) & RBV Value & H15: RQ $\rightarrow$ AT & PLS-SEM path coefficient \\
\addlinespace
F4--F6 & Explainability Value (XV) & RBV Inimitability & H16: XV $\rightarrow$ AT & PLS-SEM path coefficient \\
\addlinespace
F7--F8 & Hallucination Concerns (HC) & GenAI Risk & H17: HC $\rightarrow$ AT (negative) & PLS-SEM (F7 reverse-scored); $f^2$ effect size \\
\addlinespace
G1--G7 & Implementation Barriers & TOE + RBV & Qualitative triangulation & Thematic analysis \cite{braun2006thematic}; frequency ranking \\
\addlinespace
H1--H3 & Wearable Acceptance (WA) & TAM (Wearable) & H18: WA $\rightarrow$ PU & PLS-SEM path coefficient \\
\addlinespace
H4--H6 & Sensor Fusion Value (SF) & RBV Non-substitutable & H19: SF $\rightarrow$ MDV & PLS-SEM path coefficient \\
\addlinespace
H7--H8 & Patient Experience (PX) & TAM (Patient) & H20: PX $\rightarrow$ WA & PLS-SEM path coefficient; MGA by role (A1) \\
\end{xltabular}}
\vspace{0.5em}{\small\noindent\textit{Note.} TAM = Technology Acceptance Model; TOE = Technology--Organization--Environment; RBV = Resource-Based View; DSR = Design Science Research; PLS-SEM = Partial Least Squares Structural Equation Modelling; MGA = Multi-Group Analysis; VRIN = Valuable, Rare, Inimitable, Non-substitutable; $\beta$ = standardised path coefficient; $R^2$ = coefficient of determination; $Q^2$ = Stone--Geisser predictive relevance; $f^2$ = effect size. All hypotheses tested at $\alpha = .05$ with Bonferroni correction. Mediation effects tested using 5,000 bootstrap resamples with bias-corrected confidence intervals.}

%===============================================================================
%% ITEM SUMMARY AND RELIABILITY TARGETS
%===============================================================================
\subsection{Survey Item Summary and Psychometric Targets}
\label{subsec:survey_summary}

Table~\ref{tab:survey_item_summary} provides a summary count of items per section and the target psychometric properties for each multi-item scale.
\clearpage
{\tablestripes\tablefontsize

\noindent Table~\ref{tab:survey_item_summary} presents the survey item count and psychometric reliability targets.

\begin{xltabular}{\textwidth}{@{} L >{\centering\arraybackslash}p{1.1cm} L >{\centering\arraybackslash}p{1.6cm} >{\centering\arraybackslash}p{1.6cm} @{}}
\caption[Survey Item Count and Psychometric Reliability Targets]{Survey Item Count and Psychometric Reliability Targets by Section}\label{tab:survey_item_summary} \\
\toprule
\thead{Section} & \thead{Items} & \thead{Response Type} & \thead{Target $\alpha$} & \thead{Target AVE} \\
\midrule
\endfirsthead
\endhead
\bottomrule
\endlastfoot
A: Demographics and Professional Background & 10 & Categorical / multi-select & --- & --- \\
\addlinespace
B1--B6: Perceived Usefulness & 6 & 5-point Likert & $\geq .85$ & $\geq .55$ \\
\addlinespace
B7--B12: Perceived Ease of Use & 6 & 5-point Likert & $\geq .85$ & $\geq .55$ \\
\addlinespace
B13--B15: Behavioural Intention & 3 & 5-point Likert & $\geq .80$ & $\geq .60$ \\
\addlinespace
C1--C4: Technology Context & 4 & 5-point Likert & $\geq .80$ & $\geq .50$ \\
\addlinespace
C5--C8: Organization Context & 4 & 5-point Likert & $\geq .80$ & $\geq .50$ \\
\addlinespace
C9--C12: Environment Context & 4 & 5-point Likert & $\geq .80$ & $\geq .50$ \\
\addlinespace
D1--D3: Algorithmic Trust & 3 & 5-point Likert & $\geq .80$ & $\geq .60$ \\
\addlinespace
D4--D6: Data Governance & 3 & 5-point Likert & $\geq .80$ & $\geq .60$ \\
\addlinespace
D7--D9: Ethical Considerations & 3 & 5-point Likert & $\geq .80$ & $\geq .60$ \\
\addlinespace
D10--D12: Regulatory Compliance & 3 & 5-point Likert & $\geq .80$ & $\geq .60$ \\
\addlinespace
E1--E3: Diagnostic Performance & 3 & 5-point Likert & $\geq .80$ & $\geq .60$ \\
\addlinespace
E4--E6: Workflow Integration & 3 & 5-point Likert & $\geq .80$ & $\geq .60$ \\
\addlinespace
E7--E10: ASD-Focused Value & 4 & 5-point Likert & $\geq .80$ & $\geq .50$ \\
\addlinespace
F1--F3: RAG Report Quality & 3 & 5-point Likert & $\geq .80$ & $\geq .60$ \\
\addlinespace
F4--F6: Explainability Value & 3 & 5-point Likert & $\geq .80$ & $\geq .60$ \\
\addlinespace
F7--F8: Hallucination Concerns & 2 & 5-point Likert & $\geq .70$ & $\geq .50$ \\
\addlinespace
G1--G7: Implementation Barriers & 7 & Ranked / categorical / open-ended & --- & --- \\
\addlinespace
H1--H3: Wearable Acceptance & 3 & 5-point Likert & $\geq .80$ & $\geq .60$ \\
\addlinespace
H4--H6: Sensor Fusion Value & 3 & 5-point Likert & $\geq .80$ & $\geq .60$ \\
\addlinespace
H7--H8: Patient Experience & 2 & 5-point Likert & $\geq .70$ & $\geq .50$ \\
\midrule
\textbf{Total} & \textbf{82} & \textbf{Mixed Likert + open} & \textbf{$\geq .70$} & \textbf{$\geq .50$} \\
\end{xltabular}}
\vspace{0.5em}{\small\noindent\textit{Note.} $\alpha$ = Cronbach's alpha; AVE = Average Variance Extracted. Targets follow Hair et al.\ (2021): $\alpha \geq .70$ (acceptable), $\geq .80$ (good); AVE $\geq .50$ (adequate convergent validity). Sections A and G are not multi-item Likert scales and do not have reliability targets. Two-item scales use Spearman--Brown prophecy coefficient instead of $\alpha$.}

%===============================================================================
%% PROPOSED ANALYSIS PLAN
%===============================================================================
\subsection{Proposed Statistical Analysis Plan}
\label{subsec:survey_analysis}

\clearpage
{\tablestripes\tablefontsize

\noindent Table~\ref{tab:survey_analysis_full} presents the complete statistical analysis plan for multi-stakeholder.

\begin{xltabular}{\textwidth}{@{} p{1.0cm} p{3.4cm} p{6.5cm} p{3.2cm} @{}}
\caption[Complete Statistical Analysis Plan for Multi-Stakeholder]{Complete Statistical Analysis Plan for Multi-Stakeholder Survey}\label{tab:survey_analysis_full} \\
\toprule
\thead{\#} & \thead{Analysis Step} & \thead{Purpose} & \thead{Tool / Criterion} \\
\midrule
\endfirsthead
\endhead
\bottomrule
\endlastfoot
1 & Data screening & Missing data pattern (MCAR test), outlier detection (Mahalanobis distance), normality (skewness $< |2|$, kurtosis $< |7|$) & SPSS 29 / R; Little's MCAR test \\
\addlinespace
2 & Descriptive statistics & Sample profile by stakeholder group, response distributions, item means and standard deviations & SPSS 29 \\
\addlinespace
3 & Exploratory Factor Analysis (EFA) & Dimensionality check on pilot data ($n = 30$); confirm factor structure before CFA & SPSS; KMO $\geq .70$, Bartlett's $p < .001$ \\
\addlinespace
4 & Confirmatory Factor Analysis (CFA) & Validate measurement model; assess model fit and construct validity & SmartPLS 4; SRMR $< .08$, NFI $> .90$ \\
\addlinespace
5 & Reliability analysis & Internal consistency: Cronbach's $\alpha \geq .70$; Composite Reliability CR $\geq .70$; rho\_A $\geq .70$ & SmartPLS 4 \\
\addlinespace
6 & Convergent validity & Average Variance Extracted AVE $\geq .50$ for all reflective constructs & SmartPLS 4 \\
\addlinespace
7 & Discriminant validity & Fornell--Larcker criterion (square root of AVE $>$ inter-construct correlations); HTMT $< .85$ & SmartPLS 4 \\
\addlinespace
8 & Common method bias & Harman's single-factor test ($< 50\%$ variance on one factor); full collinearity VIF $< 3.3$ & SPSS / SmartPLS \\
\addlinespace
9 & Structural model (PLS-SEM) & Path coefficients ($\beta$), significance ($p < .05$, 5,000 bootstraps), $R^2$ ($\geq .25$ moderate), $f^2$ effect sizes & SmartPLS 4 \\
\addlinespace
10 & Mediation analysis & Bootstrapped indirect effects (5,000 resamples, 95\% bias-corrected CI) for H7 (AT mediates PU $\rightarrow$ BI) & SmartPLS 4; specific indirect effects \\
\addlinespace
11 & Moderation analysis & Interaction terms for H5 (OC moderates PU $\rightarrow$ BI), H9 (ETH moderates AT $\rightarrow$ BI) & SmartPLS 4; product indicator approach \\
\addlinespace
12 & Multi-group analysis (MGA) & Compare path coefficients across stakeholder groups (A1); parametric and permutation tests & SmartPLS 4 MGA; $p_{\text{diff}} < .05$ \\
\addlinespace
13 & Importance--Performance Map & Identify high-importance, low-performance constructs for managerial prioritization & SmartPLS 4 IPMA \\
\addlinespace
14 & Qualitative analysis (G5, G7) & Reflexive thematic analysis of open-ended responses; code--category--theme hierarchy & NVivo 14; \cite{braun2006thematic} \\
\addlinespace
15 & Triangulation & Integrate quantitative (PLS-SEM) and qualitative (thematic) findings using joint display matrix & Mixed methods joint display \\
\end{xltabular}}
\vspace{0.5em}\noindent\textit{Note.} PLS-SEM = Partial Least Squares Structural Equation Modelling; CFA = confirmatory factor analysis; EFA = exploratory factor analysis; KMO = Kaiser--Meyer--Olkin measure; SRMR = Standardized Root Mean Square Residual; NFI = Normed Fit Index; HTMT = Heterotrait-Monotrait ratio; VIF = Variance Inflation Factor; CI = confidence interval; IPMA = Importance--Performance Map Analysis; MGA = Multi-Group Analysis. All significance tests use $\alpha = .05$ with Bonferroni correction where multiple hypotheses are tested simultaneously.

%===============================================================================
%% STAKEHOLDER ROUTING MATRIX
%===============================================================================
\subsection{Stakeholder Routing Matrix}
\label{subsec:survey_routing}

Table~\ref{tab:survey_routing} specifies which survey sections are administered to each stakeholder group. All groups complete Sections A (Demographics), D (Trust and Governance), and G (Implementation Barriers). Remaining sections are routed based on role relevance to maximize response quality and minimize respondent fatigue.
\noindent Table~\ref{tab:survey_routing} stakeholder routing matrix: sections administered by respondent role.

{\tablestripes\tablefontsize
\needspace{20\baselineskip}
\begin{xltabular}{\textwidth}{@{} L c c c c c c c c r r @{}}
\caption[Stakeholder Routing Matrix]{Stakeholder Routing Matrix: Sections Administered by Respondent Role}\label{tab:survey_routing} \\
\toprule
\thead{Stakeholder Group} & \thead{A} & \thead{B} & \thead{C} & \thead{D} & \thead{E} & \thead{F} & \thead{G} & \thead{H} & \thead{Items} & \thead{Min.} \\
\midrule
\endfirsthead
\endhead
\bottomrule
\endlastfoot
Clinical practitioners & $\bullet$ & $\bullet$ & $\bullet$ & $\bullet$ & $\bullet$ & $\bullet$ & $\bullet$ & $\bullet$ & 82 & 25 \\
Healthcare administrators & $\bullet$ & $\bullet$ & $\bullet$ & $\bullet$ & -- & -- & $\bullet$ & -- & 56 & 17 \\
AI/ML engineers & $\bullet$ & $\bullet$ & -- & $\bullet$ & $\bullet$ & $\bullet$ & $\bullet$ & $\bullet$ & 70 & 21 \\
Patients and end users & $\bullet$ & -- & -- & $\bullet$ & $\bullet$ & -- & $\bullet$ & $\bullet$ & 47 & 15 \\
Regulatory officers & $\bullet$ & -- & $\bullet$ & $\bullet$ & -- & -- & $\bullet$ & -- & 41 & 14 \\
\end{xltabular}}
\vspace{0.5em}\noindent\textit{Note.} $\bullet$ = section administered; -- = section not administered (skip logic in Qualtrics/REDCap). Items = total questionnaire items administered per group. Min.\ = estimated completion time in minutes.

%===============================================================================
%% ETHICAL CONSIDERATIONS AND INFORMED CONSENT
%===============================================================================
\subsection{Ethical Considerations and Informed Consent}
\label{subsec:survey_ethics}

The following ethical safeguards apply to the administration of this survey instrument:

\begin{enumerate}
    \item \textbf{IRB Approval.} Full Institutional Review Board (IRB) approval from Golden Gate University and each participating institution is required prior to data collection. The study is classified as minimal risk (online survey, no intervention, no vulnerable population beyond standard clinical research).

    \item \textbf{Informed Consent.} The first page of the online survey presents a detailed informed consent form that describes: (a)~the purpose of the study, (b)~the voluntary nature of participation, (c)~the estimated completion time (20--25 minutes), (d)~the types of questions asked, (e)~data confidentiality and anonymization procedures, (f)~the right to withdraw at any time without penalty, and (g)~contact information for the principal investigator and IRB. Respondents must actively select ``I consent to participate'' before accessing survey items.

    \item \textbf{Anonymity and Confidentiality.} No personally identifiable information (PII) is collected. IP addresses are not logged. Institutional affiliation is collected at the category level (academic medical center, community hospital, etc.) only. All data are stored on an encrypted, access-controlled server in compliance with HIPAA and GDPR requirements.

    \item \textbf{Data Retention.} Survey responses are retained for 7 years following publication, consistent with Golden Gate University research data retention policies. After the retention period, all electronic records are permanently deleted.

    \item \textbf{Compensation.} Respondents may be offered a \$25 gift card or equivalent professional development credit as compensation for their time, subject to institutional policies. Compensation is provided regardless of survey completion to avoid coercion.

    \item \textbf{Patient Respondents.} For patient participants (Stakeholder Group~4), the informed consent form includes additional language explaining: (a)~the study is about their opinions on AI tools, not a clinical trial, (b)~their responses will not affect their care, and (c)~the survey uses simplified language (Flesch--Kincaid Grade Level $\leq 10$) to ensure accessibility.
\end{enumerate}

%===============================================================================
%% Requires prospective validation TESTING PROTOCOL
%===============================================================================
\subsection{Pilot Testing Protocol}
\label{subsec:survey_pilot}

Prior to full deployment, the survey instrument undergoes a three-stage validation process:

\begin{enumerate}
    \item \textbf{Stage 1: Expert Panel Review ($k = 5$).} Five subject-matter experts (2 neurologists, 1 healthcare informatics researcher, 1 psychometrician, 1 regulatory affairs specialist) review all items for content validity, clarity, and relevance. Items are retained if Content Validity Index (CVI) $\geq .80$ across all raters. Items with CVI $< .80$ are revised or eliminated.

    \item \textbf{Stage 2: Cognitive Interviewing ($n = 8$).} Eight participants (2 per major stakeholder group) complete the survey while thinking aloud. The interviewer probes for: (a)~comprehension difficulties, (b)~ambiguous wording, (c)~missing response options, and (d)~survey fatigue. Items are revised based on cognitive interview findings.

    \item \textbf{Stage 3: Pilot Evaluation ($n = 30$).} Thirty participants complete the survey online. Pilot data are used to assess: (a)~item distributions (ceiling/floor effects), (b)~inter-item correlations ($r \geq .30$ within constructs), (c)~exploratory factor structure, (d)~completion time, and (e)~dropout rates. Items with poor psychometric properties (item-total correlation $< .30$, cross-loading $> .40$) are revised or removed before full deployment.
\end{enumerate}

\noindent This survey instrument is designed to complement the computational validation of the RGAIG framework presented in Chapters~\ref{chap:analysis}--\ref{chap:discussion} by providing evidence from the \textit{human adoption} perspective. While the dissertation demonstrates that the framework performs technically (ASD-focused classification accuracy of 91--100\%), the survey will test whether the five stakeholder groups are \textit{willing and able} to adopt it---a necessary condition for translating computational performance into real-world clinical impact.
